\documentclass[journal]{IEEEtran}
\usepackage{fontspec}
\newfontfamily\hebrewfont[Path=/System/Library/Fonts/]{SFHebrew.ttf}
\usepackage{amsmath,amssymb,amsthm}
\usepackage{bm}
\usepackage{microtype}
\usepackage{graphicx}
\usepackage{float}
\usepackage{multicol}
\usepackage[hidelinks]{hyperref}
\setlength{\textfloatsep}{8pt plus 2pt minus 2pt}
\setlength{\floatsep}{7pt plus 2pt minus 2pt}
\setlength{\intextsep}{7pt plus 2pt minus 2pt}

\newtheorem{definition}{Definition}
\newtheorem{theorem}{Theorem}
\newtheorem{proposition}{Proposition}
\newtheorem{corollary}{Corollary}
\newtheorem{lemma}{Lemma}
\newcommand{\R}{\mathbb{R}}
\newcommand{\Tcal}{\mathcal{T}}
\newcommand{\Acal}{\mathcal{A}}
\newcommand{\Ccal}{\mathcal{C}}
\newcommand{\Pcal}{\mathcal{P}}
\newcommand{\var}{\operatorname{var}}
\newcommand{\norm}[1]{\left\lVert#1\right\rVert}

\title{Compact Ordered Nodal Variation\texorpdfstring{$^{*}$}{*}}
\author{Anonymous}
\date{}

\begin{document}
\maketitle
\begingroup
\renewcommand{\thefootnote}{\fnsymbol{footnote}}
\footnotetext[1]{The abbreviation CONV invites an implicit implementation
question: can the formal operator be compiled into a convolutional
realization? Section~XII answers that question.}
\endgroup

\begin{abstract}
Glory first be to The Holy One, Blessed be He, from whom all good things
come. We introduce Compact Ordered Nodal Variation (CONV), a
deterministic method for interpolation and matched
multiresolution analysis/synthesis on Cartesian grids. CONV combines
high-order local reconstruction, an exact factorwise variation bound, and
matched conservative multiresolution transfer. The interpolant accepts
arbitrary endpoint-aligned target lattices; dyadic analysis and integer-scale
polyphase evaluation are specializations of that continuous operator.
Within each Cartesian factor, ringing is suppressed by projecting the
high-order proposal onto a signed current fibre whose derivative-sign
variation cannot exceed that of the sampled differences.
\end{abstract}

\begin{IEEEkeywords}
Interpolation, resampling, variation diminishing, Bernstein polynomial,
Eikonal coordinates, structure tensor, lifting scheme, reversible transform,
ringing.
\end{IEEEkeywords}

\section{Introduction}

Regular-grid interpolation is commonly posed as reconstruction by a fixed
kernel. Nearest-neighbor and multilinear rules preserve range, but attenuate
high spatial frequencies; cubic convolution and windowed-sinc kernels recover
more of the passband by using signed lobes, thereby admitting new extrema near
strong transitions \cite{keys1981,lanczos1956}. Edge-directed methods instead
adapt the reconstruction to local image structure \cite{li2001,lottes2021}.
Although their direction estimates improve selected geometries, neither a
law for transported variation nor an inverse representation follows from the
direction estimate alone.

Compact Ordered Nodal Variation (CONV) is obtained by treating the observed
forward differences as signed currents on a regular Cartesian grid. Here
``ordered'' refers to their admitted sign lineage. A fourth-order interior
two-jet yields a quintic proposal, which is projected in Euclidean current
coordinates onto the observed order of variation and then integrated in the
Bernstein basis. Interpolation consequently involves three distinct objects:
the endpoint difference, its sign lineage, and the within-cell distribution
of current. Keeping these objects separate is what allows fidelity to be
improved without surrendering the variation law.

Indeed, total positivity shows that an admitted one-dimensional factor cannot
synthesize a derivative sign pair absent from the sampled ledger, while
cardinality preserves both the observed point values and the total current of
each cell. The same predictor is embedded in a dyadic transform using positive
block averages, such that cell mass is preserved and the Nyquist component is
annihilated. When the residual centered moments are retained, the resulting
analysis and synthesis maps are exact inverses. Fixed-order Cartesian
composition then carries the same construction to channel arrays.

In two dimensions, the two Cartesian factor orders are coupled by a
source-conditioned Riemannian action. The nodal action identity reduces every
source-site moment to a Kronecker value; therefore, the apparent frame bank
and nested chart family collapse exactly to the reverse factor order.
Structure-tensor anisotropy supplies the remaining scalar admission in
\([0,1]\), and the final operator is obtained as a convex blend of the two
orders without source ownership or a source-specific raster. When a global
continuation is required, the same nodal coordinate also yields the
weighted-harmonic construction of Section~XII.

Since CONV defines a continuous source-cell profile before selecting a target
grid, arbitrary \(N\to M\) synthesis merely evaluates that profile at the
endpoint-aligned target coordinates. Dyadic cell-average analysis and
integer-phase compilation are consequently specializations of one
interpolant: the former supplies a matched reversible state, whereas the
latter supplies a repeated phase bank for efficient evaluation.

\section{Relation to Existing Reconstruction}

\subsection{Sampling and raster reconstruction}

Classical sampling theory obtains exact bandlimited reconstruction from the
sinc kernel \cite{shannon1949}. In practice, finite implementations replace
sinc by a windowed kernel, of which Lanczos reconstruction is a widely used
instance \cite{lanczos1956}; cubic convolution instead trades spectral reach
for shorter support and polynomial accuracy \cite{keys1981}. Linearity and
translation invariance make the frequency response of these methods
transparent. However, a nontrivial finite low-pass approximation generally
involves signed coefficients, and its oscillatory response is consequently
unconstrained by the sign topology of the sampled differences.

Raster reconstruction developed along a complementary engineering line in
which local covariance is used to orient an edge-directed interpolant
\cite{li2001}. EASU, for example, evaluates a single-pass radial Lanczos
approximation adapted in direction and anisotropy; since the resulting signed
lobes can ring, the reference construction limits the output by the nearest
four-sample range \cite{lottes2021}. These methods make two-dimensional
structure and independent target-pixel evaluation primary design concerns.
CONV addresses the same transfer problem at an earlier representational
stage: differential current is admitted before polynomial synthesis, rather
than correcting a value after it has been synthesized.

\subsection{Shape-preserving spline interpolation}

Shape-preserving interpolation has a substantial independent history.
Monotone cubic Hermite methods constrain nodal derivatives using neighboring
secants \cite{fritsch1980,hyman1983}. Costantini constructed local
co-monotone splines of arbitrary degree from Bernstein polynomials
\cite{costantini1987}. Dougherty, Edelman, and Hyman derived derivative
constraints and modification algorithms for shape-preserving cubic and
quintic Hermite interpolation \cite{dougherty1989}. Ulrich and Watson later
gave sharp necessary-and-sufficient positivity conditions for the quartic
derivative of a monotone quintic Hermite segment \cite{ulrich1994}. MQSI uses
those conditions with local first- and second-derivative estimates and a
B-spline realization to construct a \(C^2\) monotone quintic in linear time
\cite{lux2023}.

Thus monotone quintic interpolation and its bound on derivative sign changes
are established results. In a co-monotone construction, each knot interval
receives the sign of its own secant. CONV instead retains the order of the
observed secant transitions while allowing a witnessed transition to occupy
the nearby current slots selected by \eqref{eq:boundary}; consequently, the
transition is not pinned to the sampling knot. The signed Bernstein fibre
\eqref{eq:fibre} is therefore a convex sufficient region, rather than the full
quintic monotonicity region characterized in \cite{ulrich1994}, and its
particular geometry yields Euclidean admission, endpoint-current
conservation, and the variation-diminishing proof within one finite operator.

Linear complexity alone does not make the execution structures identical.
MQSI involves derivative refinement followed by construction of a B-spline
representation \cite{lux2023}. At a fixed scale, Section~XII instead reduces
CONV to FIR proposal channels, one scalar-threshold admission, and a
polyphase synthesis matrix. The resulting bounded local primitive is the
point at which the shape-preserving lineage meets raster reconstruction.

Shape preservation on rectangular grids also predates this work. Carlson and
Fritsch gave derivative constraints and an algorithm for monotone piecewise
bicubic interpolation \cite{carlsonfritsch1989}; Costantini and Manni built
local differentiable co-monotone surfaces from variable-degree Bezier
operators \cite{costantinimanni1991}. CONV's Cartesian construction is instead
an ordered composition of the same one-dimensional current operator. The
factorwise theorem proved below is the property of that composition; it is not
a claim of first bivariate shape preservation.

\subsection{Total positivity and product representations}

Bernstein bases are nonnegative, form a partition of unity, and are totally
positive. Consequently, a Bernstein polynomial has no more sign changes than
its ordered coefficient sequence \cite{karlin1968,pinkus2010}. Variation
diminishing results also exist for tensor-product Bernstein operators and
special multivariate domains \cite{cavaretta1987,goodman1987}. These results
do not supply a general arbitrary-path variation theorem for a tensor-product
surface. We therefore state and prove the property used by grid transfer:
variation cannot be created within any admitted factor. Transverse features
formed by composition are products of recorded factor currents, not
unrecorded factor oscillations.

\subsection{Adaptive and invertible transfer}

Edge-directed interpolation estimates orientations from the values being
resampled \cite{li2001,lottes2021}. CONV instead derives its admissible
current order from the full signed differential lineage of the sampled
field. Its multidimensional Eikonal geometry likewise derives its Riemannian
action supports and directional admissions from the source current; it
accepts no external orientation as an additional input.

Harten's multiresolution algorithms already represent a hierarchy of nested
cell averages by the coarsest averages and the errors made when predicting
finer-grid averages from the next coarser level \cite{harten1995}. Cell
averages plus prediction residuals are therefore established multiresolution
coordinates. CONV uses that architecture with a different predictor: the
fine centered moment is predicted by the same admitted ordered-current
reconstruction used for interpolation, and its zero-detail synthesis obeys
the variation law proved in Section~IX.

Lifting schemes convert a prediction rule into an invertible transform by
storing its residual \cite{sweldens1998}. Direction-adaptive lifting has also
combined content-oriented image prediction with perfect reconstruction
\cite{ding2007}. CONV applies the lifting principle directly to transported
variation. Point interpolation uses the admitted Bernstein current.
Multiresolution analysis uses the cell average as its invariant state and
predicts the complementary centered moment from the same admitted profile.
The residual moment is the fine state not represented by zero-detail
synthesis.

\section{Interpolation Requirements}

The CONV construction follows from six operator requirements, which together
determine the polynomial degree, the admission rule, and the inverse
representation used below.

\subsection{Observation consistency}

For an integer refinement \(s\), the refined line has
\((N-1)s+1\) sites. Let \(D_s\) select indices \(0,s,2s,\ldots\).
Observation consistency requires
\begin{equation}
 D_s\Tcal_s=I.
 \label{eq:reqcardinal}
\end{equation}
Constants and affine coordinate fields must also be reproduced:
\begin{equation}
 \Tcal_s(\alpha i+\gamma)(r/s)=\alpha r/s+\gamma.
 \label{eq:reqaffine}
\end{equation}
Consequently, neither post-synthesis sharpening nor range correction may move
an observed sample.

\subsection{Endpoint-compatible current coordinates}

The current coordinates are chosen such that their cell mass equals the
observed endpoint difference. Thus, if \(d_h\) is the fine forward difference
and \(B_s\) sums a block of \(s\) consecutive fine differences, one obtains
\begin{equation}
 B_s d_h\Tcal_s=d_c,
 \label{eq:reqcurrent}
\end{equation}
where \(d_c y=(\delta_0,\ldots,\delta_{N-2})\). Since the endpoints are
cardinal, \eqref{eq:reqcurrent} follows by telescoping. It also gives the
representational identification used throughout the paper: internal current
mass and restricted output difference are the same observed quantity.

\subsection{Variation order}

Range preservation and absence of ringing are not identical: a function may
remain inside the global sample range while creating an alternating
derivative pair, whereas an existing extremum may move between samples
without creating such a pair. CONV therefore controls the variation topology
through the requirement
\begin{equation}
 V((\Tcal y)')\leq
 \var(\delta_0,\ldots,\delta_{N-2}).
 \label{eq:reqvariation}
\end{equation}
from which it can be seen that no new alternating rise--fall or fall--rise
pair is admitted, although an observed turning point need not remain pinned
to a grid node.

\subsection{Directness and reversibility}

Direct synthesis involves a finite local polynomial evaluation and a
deterministic finite admission step; no iterative fit or selected stopping
condition is present. Embedding the same predictor in the standard
analysis--prediction--residual lifting construction yields an analysis map
\(\mathcal L\) whose retained detail obeys
\begin{equation}
 \mathcal L^{-1}\mathcal L=I.
 \label{eq:reqinverse}
\end{equation}

\subsection{Canonical degree}

An endpoint two-jet contains the value, first derivative, and second
derivative at both ends of an interval, hence six scalar constraints. The
least polynomial degree carrying all six without an additional choice is
five; differentiation then yields five Bernstein coefficients, which form
the finite current state. The quintic degree is consequently fixed by the
compact two-jet rather than selected as a smoothness parameter.

\section{Problem and Notation}

Let \(y=(y_0,\ldots,y_{N-1})\), \(N\geq5\), be samples on a unit-spaced
one-dimensional grid. Multichannel data are handled componentwise in a
declared basis. Define
\begin{equation}
  \delta_i=y_{i+1}-y_i,\qquad 0\leq i<N-1.
  \label{eq:secant}
\end{equation}
For a finite real sequence \(z\), let \(\var(z)\) denote its number of sign
changes after zero entries are deleted. For a piecewise function \(g\), let
\(V(g)\) be the supremum of \(\var(g(x_1),\ldots,g(x_m))\) over ordered
points within open cells and one-sided knot limits, again omitting zeros.
The Bernstein basis of degree \(n\) is
\begin{equation}
  B_j^n(u)=\binom{n}{j}u^j(1-u)^{n-j},\quad 0\leq u\leq1.
  \label{eq:bernstein}
\end{equation}
For scale vector \(\bm s\), the fixed Cartesian order is part of the
operator:
\begin{equation}
  \Tcal_{\bm s}
  =\Tcal_{s_d}^{(d)}\circ\cdots\circ\Tcal_{s_1}^{(1)}.
  \label{eq:product}
\end{equation}
The reverse order is used by the corresponding analysis transform.

\subsection{Continuous query coordinates}

CONV defines a continuous piecewise-polynomial interpolant before a target
lattice is chosen. For a target line of length \(M\), endpoints are aligned
by
\begin{equation}
 x_r=\frac{r(N-1)}{M-1},\qquad 0\leq r<M.
 \label{eq:targetmap}
\end{equation}
Integer refinement is obtained when \(M=(N-1)s+1\), in which case
\eqref{eq:reqcardinal} applies exactly. Using the same coordinate map on each
Cartesian axis makes arbitrary \(N\to M\) evaluation the primary
interpolation map; the dyadic transform of Section~IX and the integer-scale
phase bank of Section~XII are, respectively, its conservative
analysis/synthesis and compiled-evaluation specializations.

\subsection{Operator factorization}

The one-dimensional method factors as
\begin{equation}
 y\xrightarrow{\mathcal J}a
 \xrightarrow{\mathcal E}\sigma
 \xrightarrow{\Pi_{\Ccal(\sigma,\delta)}}c
 \xrightarrow{\mathcal I}Q
 \xrightarrow{\mathcal B}\Tcal y.
 \label{eq:pipeline}
\end{equation}
Here \(\mathcal J\) forms the compact two-jet proposal, after which
\(\mathcal E\) orders the observed sign lineages and \(\Pi_{\Ccal}\)
performs Euclidean nearest-point admission in current coordinates. Finally,
\(\mathcal I\) integrates the admitted current and \(\mathcal B\) evaluates
the Bernstein curve. The equations below specify each arrow, such that value
observations, differential topology, and polynomial evaluation remain
separate throughout the construction.

\section{Compact Two-Jet}

The raw local model is obtained without fitted parameters. At an interior
node \(2\leq i\leq N-3\), the two-jet is
\begin{align}
 m_i&=\frac{y_{i-2}-8y_{i-1}+8y_{i+1}-y_{i+2}}{12}, \label{eq:mint}\\
 q_i&=\frac{-y_{i+2}+16y_{i+1}-30y_i+16y_{i-1}-y_{i-2}}{12}.
 \label{eq:qint}
\end{align}
Second-order one-sided closure is used at the outer two nodes:
\begin{align}
 m_0&=\frac{-3y_0+4y_1-y_2}{2},&
 m_1&=\frac{y_2-y_0}{2}, \label{eq:mbnd}\\
 q_0&=2y_0-5y_1+4y_2-y_3,&
 q_1&=y_0-2y_1+y_2, \label{eq:qbnd}
\end{align}
with the reflected formulas used at \(N-2\) and \(N-1\). Since
\eqref{eq:mint}--\eqref{eq:qbnd} are exact on affine data, the same equations
both close the boundary and preserve first-degree reproduction.

Using the two endpoint jets on cell \(i\) yields the six raw quintic Bezier
controls
\begin{align}
 P_{i0}&=y_i,&
 P_{i1}&=y_i+\frac{m_i}{5}, \nonumber\\
 P_{i2}&=y_i+\frac{2m_i}{5}+\frac{q_i}{20},&
 P_{i3}&=y_{i+1}-\frac{2m_{i+1}}{5}+\frac{q_{i+1}}{20}, \label{eq:controls}\\
 P_{i4}&=y_{i+1}-\frac{m_{i+1}}{5},&
 P_{i5}&=y_{i+1}. \nonumber
\end{align}
The corresponding raw variation current is obtained by differencing adjacent
controls,
\begin{equation}
  a_{ij}=P_{i,j+1}-P_{ij},\qquad 0\leq j<5.
  \label{eq:rawcurrent}
\end{equation}
which, by telescoping, satisfies
\(\sum_{j=0}^{4}a_{ij}=\delta_i\) identically.

\section{Ordered Lineage and Convex Current Admission}

\subsection{Lineage ledger}

The lineage component carries the order of the sign transitions observed in
the sampled differences. Each nonzero change of secant sign introduces one
boundary in the flattened current coordinate, and these locations are
admitted in increasing order,
\begin{equation}
 0=b_0<b_1<\cdots<b_L<J,
 \label{eq:causalorder}
\end{equation}
such that the next location minimizes the squared sign-conflict objective
subject to \eqref{eq:causalorder}. A later lineage can therefore neither cross
nor precede an earlier one.

Each nonzero coarse current carries its sign. A zero secant inherits the
preceding nonzero sign; a leading zero run inherits the first following
nonzero sign. If all secants vanish, the ledger is zero. Put \(I=N-1\),
\(J=5I\), and flatten five current slots per cell. Let
\(\kappa_1<\cdots<\kappa_L\) be the coarse transition knots and set
\(b_0=0\). For transition \(\ell\), define
\begin{align}
\mathcal B_\ell
 &=\{\max(b_{\ell-1}+1,5\kappa_\ell-4),\ldots,\nonumber\\
 &\hspace{25mm}\min(J-1,5\kappa_\ell+4)\},\label{eq:candidates}\\
\mathcal N_\ell
 &=\{\max(0,5\kappa_\ell-5),\ldots,
\min(J-1,5\kappa_\ell+4)\}.\label{eq:window}
\end{align}
For candidate \(b\), let \(\sigma_k(b)\) equal the coarse sign left of the
knot for \(k<b\), and the sign right of it otherwise. The admitted boundary
is
\begin{equation}
 b_\ell=\min\arg\min_{b\in\mathcal B_\ell}
 \sum_{k\in\mathcal N_\ell}\bigl[\min(0,\sigma_k(b)a_k)\bigr]^2.
 \label{eq:boundary}
\end{equation}
The outer minimum fixes ties at the smallest index. Applying
\eqref{eq:candidates}--\eqref{eq:boundary} from left to right consequently
produces a unique sign ledger \(\sigma_{ij}\in\{-1,0,1\}\). Thus the data
supply precisely the discrete current signs that interpolation must
transport, without selecting either a kernel family or a continuous tuning
parameter.

\paragraph{Exact incremental boundary cost}
Write the objective in \eqref{eq:boundary} as \(C_\ell(b)\), and let
\(s_\ell\in\{-1,1\}\) be the coarse sign immediately left of the
transition. The right sign is \(-s_\ell\). Hence
\begin{align}
 C_\ell(b)
 &=\sum_{\substack{k\in\mathcal N_\ell\\k<b}}
 a_k^2\,\mathbf1_{s_\ell a_k<0}\nonumber\\
 &\quad+\sum_{\substack{k\in\mathcal N_\ell\\k\ge b}}
 a_k^2\,\mathbf1_{-s_\ell a_k<0}.
 \label{eq:ledgercostexplicit}
\end{align}
Moving the boundary from \(b\) to \(b+1\) changes only the contribution
of slot \(b\), giving the exact recurrence
\begin{equation}
 \boxed{C_\ell(b+1)=C_\ell(b)-s_\ell a_b|a_b|.}
 \label{eq:ledgerrecurrence}
\end{equation}
Indeed, the changed contribution is
\(a_b^2(\mathbf1_{s_\ell a_b<0}-\mathbf1_{-s_\ell a_b<0})
=-s_\ell a_b|a_b|\), including \(a_b=0\). Evaluating the first
candidate cost once and then applying \eqref{eq:ledgerrecurrence} traverses
the same set \(\mathcal B_\ell\), with the same smallest-index tie rule.
Thus the recurrence changes the arithmetic realization of
\eqref{eq:boundary}, while preserving its objective, causal constraints,
and exact-arithmetic result. The finite-precision safeguard is specified
in Section~X.

\subsection{Signed mass fibre}

Once the ledger is fixed, the admissible current on each cell belongs to the
closed convex fibre
\begin{equation}
 \Ccal_i=
 \left\{c\in\R^5:
 \bm1^\mathsf Tc=\delta_i,\quad
 \sigma_{ij}c_j\geq0\ \text{for all }j
 \right\}.
 \label{eq:fibre}
\end{equation}
and the admitted current is obtained by Euclidean projection,
\begin{equation}
 c_i=\arg\min_{c\in\Ccal_i}\norm{c-a_i}_2^2.
 \label{eq:projection}
\end{equation}
Consequently, \eqref{eq:projection} is an orthogonal projection in the
canonical Euclidean coordinates of the five current increments, rather than
a componentwise clamp; whenever the raw two-jet already belongs to the signed
mass fibre, it passes unchanged. Nonemptiness follows because the ledger
contains the sign of \(\delta_i\), and strict convexity then gives a unique
solution. Since the dimension is five, this solution is obtained directly by
a scalar threshold. Indeed, a zero ledger implies constant source data and
therefore \(a_i=c_i=0\). Otherwise
\(\sigma_{ij}\in\{-1,1\}\), and, for \(\lambda\in\R\), put
\begin{equation}
 c_{ij}(\lambda)=
 \begin{cases}
 \max\{0,a_{ij}-\lambda\},&\sigma_{ij}=+1,\\
 \min\{0,a_{ij}-\lambda\},&\sigma_{ij}=-1.
 \end{cases}
 \label{eq:thresholdcurrent}
\end{equation}
The equality multiplier is consequently a solution of
\begin{equation}
 F_i(\lambda):=\sum_{j=0}^{4}c_{ij}(\lambda)=\delta_i.
 \label{eq:thresholdmass}
\end{equation}
The function \(F_i\) is continuous, nonincreasing, and piecewise affine,
with breakpoints only at \(a_{i0},\ldots,a_{i4}\). On any interval between
successive breakpoints, its active index set is
\begin{equation}
 A(\lambda)=\{j:\sigma_{ij}(a_{ij}-\lambda)>0\},
\end{equation}
and the affine mass equation yields the candidate root
\begin{equation}
 \lambda_A=\frac{\sum_{j\in A}a_{ij}-\delta_i}{|A|}.
 \label{eq:thresholdroot}
\end{equation}
Sorting the five breakpoints exposes at most six such intervals. The unique
admitted current is therefore obtained by testing
\eqref{eq:thresholdroot} against them and applying the five clips in
\eqref{eq:thresholdcurrent}.

\begin{proposition}[Exact scalar admission]
Equations \eqref{eq:thresholdcurrent}--\eqref{eq:thresholdroot} return the
Euclidean projection \eqref{eq:projection}. If \eqref{eq:thresholdmass} has
more than one multiplier, every such multiplier gives the same current.
\end{proposition}

\begin{proof}
The Karush--Kuhn--Tucker conditions for \eqref{eq:projection} give
\eqref{eq:thresholdcurrent} for the equality multiplier \(\lambda\).
Feasibility is exactly \eqref{eq:thresholdmass}. On every interval between
successive breakpoints the active set is fixed, and solving the affine mass
equation gives \eqref{eq:thresholdroot}. Nonemptiness of \(\Ccal_i\) supplies
a feasible sign for \(\delta_i\). If both ledger signs occur, then
\(F_i(\lambda)\) tends to \(+\infty\) and \(-\infty\) at the two ends of the
real line. If only the positive sign occurs, \(\delta_i\geq0\) and the range
of \(F_i\) is \([0,\infty)\); if only the negative sign occurs,
\(\delta_i\leq0\) and its range is \((-\infty,0]\). Continuity therefore
gives a root in every case. The objective is strictly convex in \(c_i\), so
its current is unique even in a zero-mass case where a multiplier interval
may remain.
\end{proof}

\begin{lemma}[Fixed-ledger contraction]
For a fixed ledger and cell mass, Euclidean admission is firmly
nonexpansive:
\begin{equation}
 \norm{\Pi_{\Ccal_i}a-\Pi_{\Ccal_i}b}_2^2
 \leq
 (\Pi_{\Ccal_i}a-\Pi_{\Ccal_i}b)^\mathsf T(a-b).
 \label{eq:firm}
\end{equation}
Consequently \(\norm{\Pi_{\Ccal_i}a-\Pi_{\Ccal_i}b}_2
\leq\norm{a-b}_2\).
\end{lemma}

\begin{proof}
Since \(\Ccal_i\) is nonempty, closed, and convex, the standard variational
inequality for orthogonal projection gives \eqref{eq:firm}; applying
Cauchy--Schwarz then yields the contraction bound.
\end{proof}

\begin{proposition}[Continuous-energy bound for the current metric]
For a coefficient error \(e\in\R^5\), let
\(E(u)=\sum_{j=0}^4 e_j\tau_j(u)\), where
\(\tau_j(u)=\sum_{r=j+1}^5 B_r^5(u)\). On the normalized cell,
\begin{equation}
 \int_0^1 |E'(u)|^2\,du
 =e^{\mathsf T}\mathsf K e\leq5\norm{e}_2^2,
 \label{eq:currentenergybound}
\end{equation}
with the derivative-energy Gram matrix
\begin{equation}
 \mathsf K_{jk}
 =25\int_0^1 B_j^4(u)B_k^4(u)\,du
 =\frac{25\binom4j\binom4k}{9\binom8{j+k}}.
 \label{eq:currentenergygram}
\end{equation}
\end{proposition}
\begin{proof}
The Bernstein derivative identity gives \(\tau_j'=5B_j^4\), proving
the equality in \eqref{eq:currentenergybound}. Each entry of
\(\mathsf K\) is nonnegative, and the partition of unity gives
\(\sum_k\mathsf K_{jk}=25\int_0^1B_j^4(u)\,du=5\).
Symmetry therefore yields
\begin{equation}
 e^{\mathsf T}(5I_5-\mathsf K)e
 =\frac12\sum_{j,k=0}^4\mathsf K_{jk}(e_j-e_k)^2\geq0.
 \label{eq:currentenergylumping}
\end{equation}
\end{proof}
Consequently, the Euclidean objective in \eqref{eq:projection}, up to the
irrelevant factor five, is the row-sum diagonal upper bound on continuous
derivative error. On a physical cell of length \(h\), both sides of
\eqref{eq:currentenergybound} acquire the factor \(h^{-1}\).
This interpretation preserves the fixed scalar-threshold admission and
does not identify it with projection in the full Gram metric. In particular,
it supplies an energy justification for the existing coordinates, rather
than a claim of universal reconstruction optimality.

\subsection{Synthesis}

Synthesis now involves a single integration of the admitted current. Starting
from \(Q_{i0}=y_i\), put
\begin{equation}
 Q_{i,j+1}=Q_{ij}+c_{ij},\qquad 0\leq j<5.
 \label{eq:integrate}
\end{equation}
The endpoint constraint in \eqref{eq:fibre} then gives
\(Q_{i5}=y_{i+1}\), and CONV cell synthesis is obtained as
\begin{equation}
 (\Tcal y)(i+u)=\sum_{j=0}^{5}Q_{ij}B_j^5(u),
 \qquad 0\leq u\leq1.
 \label{eq:synthesis}
\end{equation}
Equivalently, the entire representation is the anchored current state
\begin{equation}
 \Acal(y)=\left(y_0,\{c_{ij}\}_{i,j}\right),\qquad
 y_i=y_0+\sum_{k<i}\sum_{j=0}^{4}c_{kj}.
 \label{eq:anchored}
\end{equation}

\section{Multidimensional CONV Construction}

\subsection{Ordered factor composition}

Let \(Y^{(0)}=Y\). For axis order \(1,\ldots,d\), define
\begin{equation}
 Y^{(k)}=\Tcal_{s_k}^{(k)}Y^{(k-1)},\qquad
 \Tcal_CY=Y^{(d)}.
 \label{eq:stage}
\end{equation}
At stage \(k\), a fresh two-jet and lineage ledger are constructed on every
line parallel to axis \(k\), so that all factor currents remain represented
rather than being collapsed to a principal orientation. Since admission is
nonlinear, the factors need not commute; the order declared in
\eqref{eq:stage} is therefore part of the operator and is carried along by a
permutation of the coordinate axes. The resulting \(\Tcal_CY\) is the
Cartesian factor proposal consumed by the source-conditioned construction
below.

\subsection{Axis-permutation equivariance}

The same one-dimensional law is used on every factor. If \(P_\pi\) permutes
the array axes and the declared factor order is permuted with them, then
\begin{equation}
 \Tcal_{\pi(\bm s)}P_\pi=P_\pi\Tcal_{\bm s}.
 \label{eq:permutation}
\end{equation}
Thus \eqref{eq:permutation} gives axis-permutation equivariance when the
operator order is carried with the axes. It does not assert rotational
isotropy: because nonlinear factors need not commute, fixed
\(x\)-then-\(y\) and \(y\)-then-\(x\) compositions may differ. In two
dimensions, synthesis uses horizontal (axis 1) followed by vertical (axis
0), and analysis uses the reverse order.

% BEGIN WITNESSED EIKONAL TRANSPORT
\subsection{Source-conditioned Eikonal factor transport}
\label{sec:witnessedtransport}

Applying the one-dimensional construction in two dimensions yields two
Cartesian compositions. Let \(\Tcal_{xy}\) apply the horizontal factor and
then the vertical factor, and let \(\Tcal_{yx}\) apply them in the reverse
order. With \(\Tcal_C=\Tcal_{xy}\) as the declared Cartesian proposal, CONV
couples the two orders through one source-conditioned Riemannian action field,
without constructing a separate raster for every source.

Let the source rectangle be
\(\Omega=[0,W-1]\times[0,H-1]\), with lattice sites \(x_p\). For channel
\(c\), define
\begin{equation}
 \begin{aligned}
 g_p^{(c)}&=(D_xY_p^{(c)},D_yY_p^{(c)})^T,\\
 G_p&=\sum_{c=1}^Cg_p^{(c)}g_p^{(c)T},
 &\omega_p&=\operatorname{tr}G_p .
 \end{aligned}
 \label{eq:witnessjet}
\end{equation}
Using the same two-jet stencils and boundary closures already present in the
Cartesian factors for \(D_x,D_y\), define the dimensionless nodal quadratic
form
\begin{equation}
 Q_p=
 \begin{cases}
 I+G_p/\omega_p,&\omega_p>0,\\
 I,&\omega_p=0,
 \end{cases}
 \qquad
 M_p=\frac{Q_p}{\det(Q_p)^{1/2}}.
 \label{eq:nodalmetric}
\end{equation}
The result is SPD with \(\det M_p=1\); scalar gradient magnitude has been
removed, leaving only the measured directional shape. For \(x\) in a source
cell, let \(b_p(x)\) be its four bilinear barycentric weights and put
\begin{equation}
 \widetilde M(x)=\sum_{p\in\operatorname{cell}(x)}b_p(x)M_p,
 \qquad
 M(x)=\frac{\widetilde M(x)}{\det(\widetilde M(x))^{1/2}}.
 \label{eq:continuousmetric}
\end{equation}
Since positive barycentric combinations preserve positive definiteness,
\eqref{eq:continuousmetric} gives one continuous determinant-one SPD field on
\(\Omega\).

For each source site \(s_a=x_a\), \(1\leq a\leq A=HW\), the action is the
intrinsic Riemannian distance
\begin{equation}
 D_a(x)=\inf_{\substack{\gamma(0)=s_a,\,\gamma(1)=x\\
                        \gamma([0,1])\subseteq\Omega}}
 \int_0^1\sqrt{\dot\gamma(t)^TM(\gamma(t))\dot\gamma(t)}\,dt.
 \label{eq:sourceaction}
\end{equation}
\begin{proposition}[Well-posed constrained action]
There exist constants \(0<\mu_-\leq\mu_+<\infty\) such that
\begin{equation}
 \mu_-I\preceq M(x)\preceq\mu_+I,\qquad x\in\Omega.
 \label{eq:actionmetricbounds}
\end{equation}
Consequently, for the length metric \(d_M\) induced by
\eqref{eq:sourceaction},
\begin{equation}
 \sqrt{\mu_-}\,\norm{x-z}_2
 \leq d_M(x,z)
 \leq \sqrt{\mu_+}\,\norm{x-z}_2 .
 \label{eq:actiondistancebounds}
\end{equation}
The constrained infimum is attained for every \(x,z\in\Omega\). Hence
\(D_a(x)=d_M(s_a,x)\) is finite and continuous, vanishes only at \(s_a\),
and is single-valued even where several minimizing curves have the same
length.
\end{proposition}

\begin{proof}
Continuity and positive definiteness of \(M\) on compact \(\Omega\) give
\eqref{eq:actionmetricbounds}. The lower inequality in
\eqref{eq:actiondistancebounds} follows by bounding the metric length of
every admissible curve below by its Euclidean length; the upper inequality
follows from the straight segment in the convex rectangle. Thus \(d_M\) is
bi-Lipschitz equivalent to the Euclidean metric, so \((\Omega,d_M)\) is a
compact length space. A compact length space is geodesic, and therefore the
infimum is attained.
\end{proof}

Equivalently, the same action is obtained as the nonnegative viscosity
solution of
\begin{equation}
 \nabla D_a(x)^TM(x)^{-1}\nabla D_a(x)=1,
 \qquad D_a(s_a)=0,
 \label{eq:eikonalpde}
\end{equation}
where paths are constrained to the closed source rectangle. In practice,
anisotropic fast marching provides a direct Cartesian evaluation of
\eqref{eq:sourceaction} \cite{mirebeau2014}.
Possible nonuniqueness of a minimizing curve introduces no branch choice in
CONV, because \eqref{eq:actionpartition} depends only on its minimum length.

Using these distances, the normalized inverse-square actions are
\begin{equation}
 w_a(x)=\frac{D_a(x)^{-2}}{\sum_bD_b(x)^{-2}},
 \qquad w_a(x)\geq0,\qquad\sum_aw_a(x)=1.
 \label{eq:actionpartition}
\end{equation}
At a source site, continuous extension of the quotient assigns unit mass to
its zero-action source site and zero mass to all others. Thus
\eqref{eq:actionpartition} is Shepard inverse-distance weighting
\cite{shepard1968} in the source-derived Riemannian distance; normalization
removes the distance units and yields a direct all-source partition without
selecting an owner.

\begin{lemma}[Kronecker action reduction]
For every pair of source sites,
\begin{equation}
 w_a(x_p)=\delta_{ap}.
 \label{eq:kroneckeraction}
\end{equation}
Consequently, for every source-indexed quantity \(F_p\),
\begin{equation}
 \sum_pw_a(x_p)F_p=F_a.
 \label{eq:nodalmomentcollapse}
\end{equation}
In particular, the action-weighted support inertia about the source site is
identically zero:
\begin{equation}
 \sum_pw_a(x_p)(x_p-s_a)(x_p-s_a)^T=0.
 \label{eq:zeroinertia}
\end{equation}
\end{lemma}

\begin{proof}
Since the Riemannian distance vanishes at its own source site and is strictly
positive between distinct lattice sites, the zero-action extension of
\eqref{eq:actionpartition} gives \eqref{eq:kroneckeraction}. Substitution then
yields \eqref{eq:nodalmomentcollapse} and \eqref{eq:zeroinertia}.
\end{proof}

The lemma removes an apparent geometric degree of freedom: any frame inferred
from \eqref{eq:zeroinertia} lies necessarily in its repeated-eigenvalue
branch. With the declared Cartesian tie \(t=e_x\), \(n=e_y\), the nonzero
nested route is therefore
\begin{equation}
 \mathcal C\!\left[
   \{\mathcal C[Y_{:,i}](v)\}_{i}
 \right](u)=\Tcal_{yx}Y(u,v),
 \label{eq:collapsedroute}
\end{equation}
independently of the source index. Consequently, a source frame bank and its
nested chart construction are algebraically redundant under the action
definitions, and \eqref{eq:collapsedroute} need be evaluated only once.

What remains is to determine how strongly the reverse order is admitted.
The current tensor already supplies the required directional energy, and its
normalized vertical component gives
\begin{equation}
 \eta_a=
 \begin{cases}
 \dfrac{(G_a)_{yy}}{\operatorname{tr}G_a},&\operatorname{tr}G_a>0,\\[5pt]
 \dfrac12,&\operatorname{tr}G_a=0.
 \end{cases}
 \label{eq:witnesscompatibility}
\end{equation}
Such normalized tensor energies are standard directional measures
\cite{bigun1987}. Since \((G_a)_{yy}\geq0\) and
\((G_a)_{yy}\leq\operatorname{tr}G_a\), equation
\eqref{eq:witnesscompatibility} gives \(0\leq\eta_a\leq1\) directly, without
an eigendecomposition, threshold, or fitted gain. The zero-current value is
one half, while exchanging the Cartesian axes sends \(\eta_a\) to
\(1-\eta_a\), exactly as it exchanges the two factor orders. The scalar
reverse-order admission is consequently
\begin{equation}
 \beta(x)=\sum_aw_a(x)\eta_a,
 \qquad 0\leq\beta(x)\leq1.
 \label{eq:actionadmission}
\end{equation}
and the complete two-dimensional CONV reconstruction is
\begin{equation}
 \boxed{
 \Tcal Y(x)=
 [1-\beta(x)]\Tcal_{xy}Y(x)+\beta(x)\Tcal_{yx}Y(x) .}
 \label{eq:witnessedblend}
\end{equation}

\begin{proposition}[Source-conditioned factor convexity]
For every target point, the two coefficients in
\eqref{eq:witnessedblend} are nonnegative and sum to one. Consequently
\(\Tcal Y(x)\) belongs to the componentwise convex hull of the two Cartesian
factor-order candidates. The operator preserves constants and reproduces
affine coordinate fields whenever the constituent factors do.
\end{proposition}

\begin{proof}
From \eqref{eq:witnesscompatibility},
\(0\leq\eta_a\leq1\); using \eqref{eq:actionpartition} then gives
\(0\leq\beta\leq\sum_aw_a=1\), and the two coefficients in
\eqref{eq:witnessedblend} sum to one. Since both constituent orders return
the same constant or affine field, their convex combination returns it as
well.
\end{proof}

\begin{proposition}[Linewise blend criterion]
Let \(\gamma:[a,b]\to\Omega\) be a \(C^1\) curve and, at every regular point,
put
\begin{equation}
 A=\Tcal_{xy}Y\circ\gamma,\qquad
 B=\Tcal_{yx}Y\circ\gamma,\qquad
 \theta=\beta\circ\gamma.
\end{equation}
Then the final reconstruction obeys the exact identity
\begin{equation}
 (\Tcal Y\circ\gamma)'
 =(1-\theta)A'+\theta B'+\theta'(B-A).
 \label{eq:blendderivative}
\end{equation}
Suppose \(sA'\geq0\) and \(sB'\geq0\) for a fixed
\(s\in\{-1,1\}\). The blend has \(s(\Tcal Y\circ\gamma)'\geq0\) exactly when
\begin{equation}
 s\theta'(B-A)
 \geq-\{(1-\theta)sA'+\theta sB'\}.
 \label{eq:blendcriterion}
\end{equation}
In particular, the readily checked bound
\begin{equation}
 |\theta'|\,|B-A|
 \leq(1-\theta)sA'+\theta sB'
 \label{eq:blendsufficient}
\end{equation}
is sufficient.
\end{proposition}

\begin{proof}
Applying the product rule to \eqref{eq:witnessedblend} gives
\eqref{eq:blendderivative}. Multiplication by \(s\) and rearrangement then
yields \eqref{eq:blendcriterion}, while \eqref{eq:blendsufficient} bounds its
left-hand term from below.
\end{proof}

\begin{proposition}[Pointwise convexity is not variation diminution]
Nonnegative pointwise blending of two monotone functions with common
endpoints does not by itself preserve monotonicity.
\end{proposition}

\begin{proof}
On \([0,1]\), take
\begin{equation}
 A(t)=t,\qquad B(t)=t+t(1-t),\qquad
 \theta(t)=\frac{1+\sin(4\pi t)}{2}.
\end{equation}
Both \(A\) and \(B\) are nondecreasing and share their endpoint values.
For \(T=(1-\theta)A+\theta B\), equation
\eqref{eq:blendderivative} gives
\begin{align}
 T'(0)&=\frac32>0,&
 T'(\tfrac34)&=\frac34-\frac{3\pi}{8}<0,\nonumber\\
 T'(1)&=\frac12>0.&&
\end{align}
Continuity therefore forces at least two derivative-sign transitions. Hence
the convex-hull statement alone cannot prove a sign-pair bound.
\end{proof}

The factorwise ringing-incapability theorem applies unconditionally to each
Cartesian composition in its declared order. The final spatially varying
blend adds no amplitude outside the componentwise interval spanned by those
two admitted values; its linewise sign topology is governed by
\eqref{eq:blendcriterion}, not by convexity alone.
% END WITNESSED EIKONAL TRANSPORT

\subsection{Componentwise multichannel interpolation}

For \(Y:\mathbb Z^d\rightarrow\R^C\), the coordinate map and factor order are
shared, while two-jets, ledgers, and fibre projections are evaluated for each
component in a declared basis. The Cartesian factor proposals are therefore
componentwise. In two dimensions, \eqref{eq:witnessjet} sums the channel
gradient outer products and supplies one shared metric and scalar admission;
no channel selects a separate factor order. Component currents remain
individually admitted inside that shared geometry. The factorwise guarantees
hold per component and are preserved by independent affine changes of channel
units; they do not assert covariance under an arbitrary basis rotation that
mixes channels.

\subsection{Compact dependence and causal state}

On interior cell \(i\), the raw current depends on the six samples from
\(y_{i-2}\) through \(y_{i+3}\). Bernstein evaluation uses only that cell's
five admitted coefficients. The lineage is formed in one causal pass. Over
an exact zero-secant run, the pass carries the last nonzero sign; a leading
run receives the first following sign. The numerical jet and polynomial are
compact. Only the sign state can traverse a plateau, preventing that plateau
from creating a false pair of fronts.

\subsection{Intercell regularity}

Every cell polynomial is smooth in its open interval and the assembled
interpolant is continuous. Its one-sided derivatives at an interior node are
\begin{align}
 (\Tcal y)'(i^-)&=5c_{i-1,4},\nonumber\\
 (\Tcal y)'(i^+)&=5c_{i,0},\label{eq:firsttrace}\\
 (\Tcal y)''(i^-)&=20(c_{i-1,4}-c_{i-1,3}),\nonumber\\
 (\Tcal y)''(i^+)&=20(c_{i,1}-c_{i,0}).
 \label{eq:secondtrace}
\end{align}
Consequently, CONV is \(C^1\) at node \(i\) exactly when
\(c_{i-1,4}=c_{i,0}\), and it is \(C^2\) there when the corresponding
second condition in \eqref{eq:secondtrace} also agrees. Since admission
imposes neither equality, a witnessed change of differential state is
retained rather than concealed merely to enforce intercell smoothness.

\section{Analytical Properties}

\begin{proposition}[Cardinality and commutation]
For every input \(y\), \((\Tcal y)(i)=y_i\). For every scalar
\(\alpha\) and constant \(\gamma\),
\[
 \Tcal(\alpha y+\gamma)=\alpha\Tcal(y)+\gamma.
\]
\end{proposition}

\begin{proof}
Since \(Q_{i0}=y_i\) and \(Q_{i5}=y_{i+1}\), cardinality follows
immediately. Translation leaves every current unchanged, while scaling
multiplies the raw current, fibre, projection, and integrated controls by
\(\alpha\); when \(\alpha<0\), the ledger reverses with the current. The
deterministic tie rules are invariant under both operations, which proves the
stated commutation.
\end{proof}

\subsection{Reproduction and ringing incapability}

\begin{proposition}[Affine reproduction]
If \(y_i=\alpha i+\gamma\), then
\[
  (\Tcal y)(x)=\alpha x+\gamma
\]
throughout the sampled interval.
\end{proposition}

\begin{proof}
The finite differences, including the one-sided closures, give
\(m_i=\alpha\) and \(q_i=0\), and therefore
\(a_{ij}=\alpha/5\). Since this current already lies in its signed fibre,
projection leaves it unchanged; \eqref{eq:synthesis} is consequently the
degree-elevated affine polynomial.
\end{proof}

\begin{corollary}[Cartesian consistency]
The product operator \eqref{eq:stage} is cardinal on the nested source
lattice and reproduces every affine coordinate field
\[
 f(x)=\gamma+\sum_{k=1}^d\alpha_kx_k.
\]
\end{corollary}

\begin{proof}
Each factor fixes its source nodes and exactly reproduces the affine
restriction of \(f\) to every line parallel to that factor. Induction over
\eqref{eq:stage} gives both statements.
\end{proof}

\begin{definition}
A \emph{factor ringing pair} is a pair of consecutive sign transitions of
the derivative synthesized inside one Cartesian factor that has no
corresponding pair in that factor's admitted current ledger.
\end{definition}

\begin{theorem}[Factorwise ringing incapability]
CONV synthesis cannot create a factor ringing pair. More
precisely, with \(c^\flat\) denoting the flattened current ledger,
\begin{equation}
 V((\Tcal y)')\leq\var(c^\flat)\leq
 \var(\delta_0,\ldots,\delta_{N-2}).
\label{eq:vardim}
\end{equation}
\end{theorem}

\begin{proof}
Differentiating \eqref{eq:synthesis} gives
\begin{equation}
 (\Tcal y)'(i+u)=5\sum_{j=0}^{4}c_{ij}B_j^4(u).
 \label{eq:derivative}
\end{equation}
Using total positivity of the Bernstein basis bounds the sign variation of
\eqref{eq:derivative} by that of its ordered coefficients. At a knot,
nonzero one-sided derivative signs come from the terminal current
coefficients of the incident cells. Those coefficients are consecutive in
\(c^\flat\). The cell bounds therefore give
\(V((\Tcal y)')\leq\var(c^\flat)\). Since admission places one ordered
boundary for each nonzero coarse sign transition, projection cannot
introduce a sign outside the ledger and can delete a transition only by
zeroing its incident coefficients. Therefore
\(\var(c^\flat)\leq\var(\delta)\), which completes
\eqref{eq:vardim}.
\end{proof}

\begin{corollary}
If the admitted ledger is nonnegative or nonpositive on a cell, synthesis is
monotone on that cell and remains within its endpoint range.
\end{corollary}

The theorem is deliberately factorwise. A multivariate tensor-product
surface may have transverse critical structure not visible on a single
coordinate line. Ordered Cartesian composition transports a separate ledger
at each stage. Such transverse crossings are formed from recorded factor
currents; they are not factor ringing and are not excluded.

\subsection{Consistency order}

\begin{proposition}[Uncontracted smooth-cell accuracy]
Restore physical mesh width \(h\), and let \(f\in C^6\). On an interior
cell \(2\leq i\leq N-4\) whose raw current is already admissible,
\begin{equation}
 \sup_{0\leq u\leq1}
 |(\Tcal f_h)(x_i+uh)-f(x_i+uh)|=O(h^5).
 \label{eq:orderinterior}
\end{equation}
On the outer and near-outer cells the stated one-sided closure gives
\(O(h^3)\).
\end{proposition}

\begin{proof}
Taylor expansion of \eqref{eq:mint} and \eqref{eq:qint} gives
\[
 m_i=hf'(x_i)+O(h^5),\qquad
 q_i=h^2f''(x_i)+O(h^6).
\]
Exact endpoint two-jets determine the quintic Hermite interpolant with
\(O(h^6)\) remainder. The first-derivative proposal therefore limits the
combined interior order to \(O(h^5)\). At the boundary,
\eqref{eq:mbnd} has scaled error \(O(h^3)\), while
\eqref{eq:qbnd} has scaled error \(O(h^4)\). If the raw current is
admissible, projection is the identity and these orders pass unchanged to
synthesis.
\end{proof}

Away from a true critical point, a sufficiently resolved smooth monotone
profile has constant secant sign and its raw current eventually lies in the
same signed fibre. At a transition, CONV retains the global absolute bound
below rather than assuming that admission is inactive.

\subsection{Absolute stability and error}

Let
\[
 |y|_{1,h}=\max_i\frac{|y_{i+1}-y_i|}{h}.
\]
Set
\[
\Delta_k=\frac{y_{k+1}-y_k}{|y|_{1,h}h}.
\]
Substitution in \eqref{eq:mint}--\eqref{eq:rawcurrent} gives
\(a_i=A_t\Delta\,|y|_{1,h}h\). Here \(t\) denotes an outer, near-outer,
or interior stencil. The exact squared row-sum certificates are
\begin{equation}
 \sum_j\left(\sum_k |(A_t)_{jk}|\right)^2
 \in\left\{\frac{49}{100},\frac{81}{100},\frac{37}{45}\right\}.
 \label{eq:certificate}
\end{equation}
The right boundary gives the reflected matrices. Since
\(\norm{\Delta}_\infty\leq1\), the largest value in
\eqref{eq:certificate} is \(37/45=185/225\), and therefore
\begin{equation}
 \norm{a_i}_2\leq\frac{\sqrt{185}}{15}|y|_{1,h}h.
 \label{eq:rawbound}
\end{equation}
The fibre contains a vector with one entry equal to \(\delta_i\) and all
others zero. Metric projection and the triangle inequality therefore give
\begin{equation}
 \norm{c_i}_2
 \leq\left(1+\frac{2\sqrt{185}}{15}\right)|y|_{1,h}h.
 \label{eq:currentbound}
\end{equation}
Because Bernstein synthesis is a convex combination of integrated controls,
\begin{equation}
 |(\Tcal y)(i+u)-y_i|
 \leq\beta |y|_{1,h}h,\qquad
\beta=\sqrt5\left(1+\frac{2\sqrt{185}}{15}\right).
\label{eq:localbound}
\end{equation}

\begin{proposition}[Amplitude excursion]
Let \(m=\min_i y_i\), \(M=\max_i y_i\), and \(R=M-m\). Define
\begin{equation}
 \operatorname{exc}(z;[m,M])=
 \max\left\{(\sup_xz-M)_+,(m-\inf_xz)_+\right\}.
 \label{eq:excursiondef}
\end{equation}
For one factor,
\begin{equation}
 \operatorname{exc}(\Tcal y;[m,M])\leq\beta R.
 \label{eq:excursionone}
\end{equation}
On every cell whose ledger has one sign, the sharper value is zero relative
to that cell's endpoint range. Hence amplitude excursion can occur only in a
cell carrying an admitted sign transition. After \(d\) ordered Cartesian
factors,
\begin{equation}
 \operatorname{exc}(\Tcal_{\bm s}Y;[m,M])
 \leq\frac{(1+2\beta)^d-1}{2}R.
 \label{eq:excursiond}
\end{equation}
These bounds are independent of the number of source or target samples.
\end{proposition}

\begin{proof}
For one factor, \( |y|_{1,h}h=\max_i|\delta_i|\leq R\), so
\eqref{eq:localbound} places every synthesized value in
\([m-\beta R,M+\beta R]\), proving \eqref{eq:excursionone}. A one-sign ledger
makes \eqref{eq:derivative} one-signed by nonnegativity of the Bernstein
basis, so the cell is monotone between its observed endpoints. For the
Cartesian bound, put \(c=(m+M)/2\). Translation equivariance and
\eqref{eq:bibo} give
\(\lVert\Tcal_{\bm s}(Y-c)\rVert_\infty
 \leq(1+2\beta)^dR/2\). Subtracting the original half-range \(R/2\) yields
\eqref{eq:excursiond}.
\end{proof}

\begin{proposition}[BIBO stability]
On one factor,
\begin{equation}
 \norm{\Tcal y}_\infty\leq(1+2\beta)\norm{y}_\infty.
 \label{eq:bibo}
\end{equation}
After \(d\) Cartesian factors the bound is
\((1+2\beta)^d\norm{Y}_\infty\).
\end{proposition}

\begin{proof}
\(|\delta_i|\leq2\norm{y}_\infty\). Apply
\eqref{eq:localbound} with \(h=1\), then compose the one-factor inequality.
\end{proof}

\begin{theorem}[Grid-size-independent bound]
Let \(f\) be \(L\)-Lipschitz and let its samples be transferred through
\(d\) ordered Cartesian factors of mesh width at most \(h\). Then
\begin{equation}
 |\widehat f(x)-f(x)|\leq E_dLh,\qquad
 E_d=(1+\beta)\frac{(1+2\beta)^d-1}{2\beta},
 \label{eq:errorbound}
\end{equation}
and
\begin{equation}
 \operatorname{MSE}(\widehat f,f)\leq E_d^2L^2h^2.
 \label{eq:msebound}
\end{equation}
The constants do not depend on lattice cardinality or the number of queried
points.
\end{theorem}

\begin{proof}
From \eqref{eq:localbound}, one obtains the one-factor amplification. Thus,
if \(E_kLh\) bounds the error after \(k\) factors, the next factor obeys
\[
 E_{k+1}=(1+2\beta)E_k+1+\beta,\qquad E_0=0.
\]
Solving this recurrence yields \eqref{eq:errorbound}; squaring and averaging
then yields \eqref{eq:msebound}.
\end{proof}

Within a chamber having fixed lineage boundaries and active clip sets, the
operator is affine and hence Lipschitz. Across chamber boundaries,
\eqref{eq:currentbound}--\eqref{eq:msebound} provide the global absolute
bound without assuming differentiability of the discrete selection map.

% BEGIN CONSERVATIVE MULTIRESOLUTION CORE
\section{Restriction and Matched Multiresolution Geometry}

\subsection{Point-current restriction}

Point interpolation and raster analysis involve different restriction
functionals. Let \(d_h\) be the forward difference on an interpolated line,
and let \(B_s\) sum each block of \(s\) consecutive fine differences. Using
cardinality gives
\begin{align}
 (B_s d_h\Tcal_s y)_i
 &=\sum_{r=0}^{s-1}
 \bigl[(\Tcal_s y)_{is+r+1}-(\Tcal_s y)_{is+r}\bigr]\nonumber\\
 &=y_{i+1}-y_i.
 \label{eq:telescoping}
\end{align}
Consequently,
\begin{equation}
 D_s\Tcal_s=I,\qquad B_s d_h\Tcal_s=d_c,
 \label{eq:restriction}
\end{equation}
where \(D_s\) selects source sites and \(d_c\) is the source forward
difference. Thus \eqref{eq:restriction} identifies the conserved mass of the
interpolating current; site selection itself is not used as an antialiased
raster downscale.

\subsection{Conservative scale coordinates}

The centered-moment construction below uses the fixed three-point
Gauss--Legendre nodes and weights
\begin{equation}
 \bm\xi=(-\sqrt{3/5},0,\sqrt{3/5}),\qquad
 \bm w=(5/9,8/9,5/9).
 \label{eq:gauss3}
\end{equation}

At a multiresolution interface, raster entries are interpreted as cell
averages. Accordingly, for an even fine line
\(x=(x_0,\ldots,x_{2N-1})\), define the parent average and centered child
moment by
\begin{align}
 z_i&=\frac{x_{2i}+x_{2i+1}}{2},
 \label{eq:scaleaverage}\\
 \mu_i^\star&=\frac{x_{2i+1}-x_{2i}}{2}.
 \label{eq:actualmoment}
\end{align}
The resulting analysis low-pass is positive and normalized, with transfer
function
\begin{equation}
 H(\omega)=\frac{1+e^{-\mathrm i\omega}}{2}
 =e^{-\mathrm i\omega/2}\cos(\omega/2),
 \label{eq:nyquistresponse}
\end{equation}
which yields \(H(\pi)=0\) exactly.

The centered-moment proposal is obtained from the already defined CONV
profile, rather than from a second interpolator. Using the Gauss--Legendre
nodes and weights in \eqref{eq:gauss3}, for \(1\leq i<N-1\), put
\begin{align}
 \bar z_i^-&=\frac12\sum_{k=0}^{2}w_k
 (\Tcal z)(i-\tfrac14+\tfrac14\xi_k),\nonumber\\
 \bar z_i^+&=\frac12\sum_{k=0}^{2}w_k
 (\Tcal z)(i+\tfrac14+\tfrac14\xi_k),
 \label{eq:halfcells}\\
 p_i&=\frac{\bar z_i^+-\bar z_i^-}{2}.
 \label{eq:momentproposal}
\end{align}
Since \(\Tcal z\) is quintic on each half-cell, three-point Gauss--Legendre
quadrature is exact. At the two outer cells, the fixed closure is
\begin{equation}
 p_0=m_0/4,\qquad p_{N-1}=m_{N-1}/4,
 \label{eq:momentboundary}
\end{equation}
using the one-sided jet in \eqref{eq:mbnd}; this proposal is exact for affine
cell averages, for which \(p_i=(z_{i+1}-z_i)/4\).

\subsection{Admissible moment polytope}

Put \(\Delta_i=z_{i+1}-z_i\), fill zero signs by the ledger rule of
Section~V, and denote the resulting interval signs by \(s_i\). When every
\(\Delta_i\) vanishes, set \(\Pcal(z,t)=\{0\}\) and
\(\widehat\mu=0\). Otherwise the filled signs are nonzero, and each cell is
assigned a sign \(t_i\), with \(t_0=s_0\),
\(t_{N-1}=s_{N-2}\), and, in the interior,
\begin{equation}
 t_i=\begin{cases}
 s_{i-1},&s_{i-1}=s_i,\\
 \arg\max_{s\in\{s_{i-1},s_i\}}sp_i,&s_{i-1}\ne s_i,
 \end{cases}
 \label{eq:momentsign}
\end{equation}
with the left sign selected on a tie. The second line therefore locates an
existing turning point in one of its two adjacent child cells.

When \(\mu\) is synthesized with zero detail, the fine currents are
\begin{equation}
 g_{2i}=2\mu_i,\qquad
 g_{2i+1}=\Delta_i-\mu_i-\mu_{i+1}.
 \label{eq:finecurrentmoment}
\end{equation}
Consequently, for fixed \(t\), the complete admissible set is
\begin{equation}
 \Pcal(z,t)=\left\{\mu:
 t_i\mu_i\geq0,\quad
 s_i(\Delta_i-\mu_i-\mu_{i+1})\geq0
 \right\}.
 \label{eq:momentpolytope}
\end{equation}
\begin{lemma}[Compact moment fibre]
For every nonconstant \(z\) and every assignment \(t\) produced by
\eqref{eq:momentsign}, \(\Pcal(z,t)\) is a nonempty compact convex polytope.
\end{lemma}

\begin{proof}
The set is a finite intersection of closed half-spaces and contains zero, so
only an unbounded ray remains to be excluded. Let \(h\) belong to its
recession cone and put \(v_i=t_ih_i\geq0\). The homogeneous face inequalities
are
\begin{equation}
 (s_it_i)v_i+(s_it_{i+1})v_{i+1}\leq0.
 \label{eq:momentrecession}
\end{equation}
Suppose \(v_k>0\) for the least such \(k\). If \(k=0\), then \(t_k=s_k\).
If \(k=N-1\), the preceding inequality contradicts \(v_k>0\) because
\(t_{N-1}=s_{N-2}\). If \(0<k<N-1\), the preceding inequality and
\(v_{k-1}=0\) either contradict \(v_k>0\), or force
\(t_k=-s_{k-1}=s_k\). Thus every surviving case has \(k<N-1\) and
\(t_k=s_k\). Equation~\eqref{eq:momentrecession} on face \(k\) then forces
\(t_{k+1}=-s_k\) and \(v_{k+1}\geq v_k\). By
\eqref{eq:momentsign}, this is possible only when
\(s_{k+1}=-s_k\) and \(t_{k+1}=s_{k+1}\). Repetition propagates a positive
nondecreasing sequence to the right endpoint, where
\(t_{N-1}=s_{N-2}\) contradicts the required opposite sign. Hence the
recession cone is \(\{0\}\), and the closed polyhedron is bounded.
\end{proof}

On an increasing run, the constraints reduce to
\begin{equation}
 \mu_i\geq0,\qquad \mu_i+\mu_{i+1}\leq\Delta_i.
 \label{eq:monotonemoment}
\end{equation}

The finite admission contains no limiter constant. First form
\begin{equation}
p_i^+=t_i\max(t_ip_i,0).
 \label{eq:momentorthant}
\end{equation}
Let \(\rho(i)\) denote the maximal sign lineage assigned to cell \(i\), and
let \(\rho(\ell)\) denote the lineage of face \(\ell\). Its proposed
consumption is
\begin{equation}
 C_\ell=\sum_{\substack{j\in\{\ell,\ell+1\}\\
                        \rho(j)=\rho(\ell)}}|p_j^+|.
 \label{eq:momentconsumption}
\end{equation}
and, for each lineage, retain the scale
\begin{equation}
\alpha_\rho=\min\left(
 1,\min_{\substack{\ell:\rho(\ell)=\rho\\C_\ell>0}}
 \frac{|\Delta_\ell|}{C_\ell}\right).
 \label{eq:momentrayscale}
\end{equation}
where the inner minimum is \(+\infty\) when the lineage has no positive
consumption. Oppositely oriented moments at a transition relax rather than
consume the face mass and therefore do not enter \(C_\ell\). The result is the
largest independent lineage-ray scale satisfying every incident capacity
bound, and it is denoted by
\begin{equation}
 \widehat\mu_i=\alpha_{\rho(i)}p_i^+=(\Acal_zp)_i.
 \label{eq:momentadmission}
\end{equation}

\begin{lemma}[Admitted-moment bound]
The admission \eqref{eq:momentadmission} satisfies
\begin{equation}
 \norm{\widehat\mu(z)}_\infty\leq
 \max_\ell|\Delta_\ell|\leq2\norm z_\infty.
 \label{eq:admittedmomentbound}
\end{equation}
\end{lemma}

\begin{proof}
Every cell \(i\) has an incident face \(\ell\) in its assigned lineage, as
forced at the endpoints and by the definition of \(t_i\) in the interior.
Consequently \(C_\ell\geq|p_i^+|\). If \(C_\ell=0\), then
\(\widehat\mu_i=0\). Otherwise
\eqref{eq:momentrayscale} gives
\[
 |\widehat\mu_i|
 =\alpha_{\rho(i)}|p_i^+|
 \leq|\Delta_\ell|.
\]
Finally, \(|z_{\ell+1}-z_\ell|\leq2\norm z_\infty\), which proves the second
inequality.
\end{proof}

\subsection{Exact conservative lifting pair}

The detail coordinate is obtained as the child moment not represented by the
transported coarse variation,
\begin{equation}
 r_i=\mu_i^\star-\widehat\mu_i(z).
 \label{eq:analysisdetail}
\end{equation}
so analysis is \(\mathcal Lx=(z,r)\), while synthesis is
\begin{align}
 x_{2i}&=z_i-\widehat\mu_i(z)-r_i,
 \label{eq:synthleft}\\
 x_{2i+1}&=z_i+\widehat\mu_i(z)+r_i.
 \label{eq:synthright}
\end{align}

\begin{theorem}[Conservative CONV lifting]
For every fine line of length \(2N\), \(N\geq5\), in exact arithmetic,
\begin{equation}
 \mathcal L^{-1}\mathcal Lx=x,
 \qquad R\mathcal L^{-1}(z,r)=z,
 \label{eq:fullcycle}
\end{equation}
where \(R\) is the positive restriction \eqref{eq:scaleaverage}. If \(r=0\),
the synthesized current has no more nonzero sign changes than \(Dz\).
\end{theorem}

\begin{proof}
Substituting \eqref{eq:analysisdetail} into
\eqref{eq:synthleft}--\eqref{eq:synthright} returns
\(z_i\mp\mu_i^\star=x_{2i},x_{2i+1}\). Their average is \(z_i\), proving
both identities. By \eqref{eq:momentpolytope}, the nonzero fine-current signs
are a subsequence of
\[
 t_0,s_0,t_1,s_1,\ldots,s_{N-2},t_{N-1}.
\]
Within a sign run the entries agree. At a transition, \(t_i\) equals one of
its two adjacent signs; that triple therefore contains one transition, not
two. Deleting zero currents cannot increase the count.
\end{proof}

The transform consequently distinguishes exact recovery from
discarded-detail reconstruction. Retaining \(r\) recovers every fine cell,
whereas setting \(r=0\) retains cell mass and transported variation while
omitting the unpredicted centered moment. For two detail states on the same
coarse raster, one obtains
\begin{equation}
 \norm{\mathcal L^{-1}(z,r)-\mathcal L^{-1}(z,r')}_\infty
 =\norm{r-r'}_\infty.
 \label{eq:detailstability}
\end{equation}

\begin{proposition}[Multilevel max-norm stability]
The positive restriction satisfies \(\norm R_\infty\leq1\). A \(J\)-level
synthesis obeys
\begin{equation}
 \norm x_\infty\leq3^J\norm{z^{(J)}}_\infty
 +\sum_{j=1}^{J}3^{j-1}\norm{r^{(j)}}_\infty.
 \label{eq:multilevelstability}
\end{equation}
The constant depends on the declared depth, not the line length.
\end{proposition}

\begin{proof}
Since the two restriction weights are positive and sum to one,
\(\norm R_\infty\leq1\). At one synthesis level, using
\eqref{eq:synthleft}--\eqref{eq:synthright} together with
\eqref{eq:admittedmomentbound} gives
\[
 \norm{x^{(j-1)}}_\infty
 \leq3\norm{x^{(j)}}_\infty+\norm{r^{(j)}}_\infty.
\]
Iterating this recurrence from \(x^{(J)}=z^{(J)}\) yields
\eqref{eq:multilevelstability}; induction on the one-level theorem likewise
shows that ordered sign variation is nonincreasing at every zero-detail
level.
\end{proof}

\subsection{Two-dimensional moment geometry}

For a fine \(2\times2\) block, introduce one invariant mean and three
centered moments,
\begin{align}
 z&=\tfrac14(x_{00}+x_{01}+x_{10}+x_{11}),\nonumber\\
 \mu_x&=\tfrac14(-x_{00}+x_{01}-x_{10}+x_{11}),\nonumber\\
 \mu_y&=\tfrac14(-x_{00}-x_{01}+x_{10}+x_{11}),\nonumber\\
 \mu_{xy}&=\tfrac14(x_{00}-x_{01}-x_{10}+x_{11}).
 \label{eq:blockmoments}
\end{align}
from which the inverse is obtained as
\begin{align}
 x_{00}&=z-\mu_x-\mu_y+\mu_{xy},&
 x_{01}&=z+\mu_x-\mu_y-\mu_{xy},\nonumber\\
 x_{10}&=z-\mu_x+\mu_y-\mu_{xy},&
 x_{11}&=z+\mu_x+\mu_y+\mu_{xy}.
 \label{eq:blockinverse}
\end{align}
Thus the four child cells are replaced without redundancy. The horizontal
and vertical proposals are obtained by applying \eqref{eq:halfcells} along
their corresponding factors, and applying the same half-cell functional to
the horizontal proposal along the vertical factor yields the mixed proposal.

Inside a block, the horizontal face currents are
\begin{equation}
 2(\mu_x-\mu_{xy}),\qquad2(\mu_x+\mu_{xy}),
 \label{eq:internalhorizontal}
\end{equation}
and the vertical face currents are
\begin{equation}
 2(\mu_y-\mu_{xy}),\qquad2(\mu_y+\mu_{xy}).
 \label{eq:internalvertical}
\end{equation}
For horizontally adjacent blocks \(A,B\), put \(\Delta_x=z_B-z_A\). Their
upper and lower cross-face currents are
\begin{align}
 h^-_{AB}&=\Delta_x-\mu_{x,A}-\mu_{x,B}
 +\mu_{y,A}-\mu_{y,B}+\mu_{xy,A}+\mu_{xy,B},\nonumber\\
 h^+_{AB}&=\Delta_x-\mu_{x,A}-\mu_{x,B}
 -\mu_{y,A}+\mu_{y,B}-\mu_{xy,A}-\mu_{xy,B}.
 \label{eq:horizontalfaces}
\end{align}
For vertically adjacent blocks \(A,C\), put \(\Delta_y=z_C-z_A\). Their
left and right currents are
\begin{align}
 v^-_{AC}&=\Delta_y+\mu_{x,A}-\mu_{x,C}
 -\mu_{y,A}-\mu_{y,C}+\mu_{xy,A}+\mu_{xy,C},\nonumber\\
 v^+_{AC}&=\Delta_y-\mu_{x,A}+\mu_{x,C}
 -\mu_{y,A}-\mu_{y,C}-\mu_{xy,A}-\mu_{xy,C}.
 \label{eq:verticalfaces}
\end{align}

Let \(t_x,t_y\) be block signs assigned by \eqref{eq:momentsign} on each
coarse row and column, and let \(s_x,s_y\) be the corresponding coarse face
signs. The face-current admissible set is
\begin{equation}
\begin{aligned}
 t_x(\mu_x\pm\mu_{xy})&\geq0,&s_xh^\pm&\geq0,\\
 t_y(\mu_y\pm\mu_{xy})&\geq0,&s_yv^\pm&\geq0.
\end{aligned}
 \label{eq:facepolytope}
\end{equation}
This admissible set is a finite intersection of affine half-spaces containing
the zero moment field. Admission first enforces the internal cone
\begin{equation}
 t_x\mu_x\geq|\mu_{xy}|,\qquad
 t_y\mu_y\geq|\mu_{xy}|.
 \label{eq:internalcone}
\end{equation}
For a cross face written \(g=\Delta+r_A+r_B\), only opposed proposal mass
consumes the witnessed capacity. Define
\begin{align}
 w_A&=\max(-sr_A,0),&w_B&=\max(-sr_B,0),\nonumber\\
 \beta&=\min\left(1,\frac{|\Delta|}{w_A+w_B}\right),
 \label{eq:facecapacity}
\end{align}
with \(\beta=1\) when the denominator is zero. Both incident block scales are
bounded by \(\beta\), and each block retains the minimum bound from all of its
faces. Since helpful components are not required for feasibility, a later
reduction on another face cannot violate an admitted inequality.
Consequently, \eqref{eq:internalcone}--\eqref{eq:facecapacity} yield a finite
local box inside \eqref{eq:facepolytope}.

Every pair of fine rows has a horizontal sign sequence interleaving
\(t_x,s_x,t_x,s_x,\ldots\); every pair of fine columns has the analogous
vertical sequence. The one-dimensional sign-count proof applies to each.
The two-dimensional restriction response is
\begin{equation}
 H(\omega_x,\omega_y)=e^{-\mathrm i(\omega_x+\omega_y)/2}
 \cos(\omega_x/2)\cos(\omega_y/2),
 \label{eq:nyquist2d}
\end{equation}
which annihilates both axial Nyquist lines and their checkerboard
intersection. Exact reconstruction is then obtained by restoring the three
residual moment fields in \eqref{eq:blockmoments}.

\subsection{Multilevel state}

Repeated analysis produces the state
\begin{equation}
 \mathcal L^{(J)}x=
 \left(z^{(J)},r^{(J)},r^{(J-1)},\ldots,r^{(1)}\right).
 \label{eq:multilevel}
\end{equation}
In one dimension, \(r^{(j)}\) is a centered-moment field; in two dimensions,
it contains the horizontal, vertical, and mixed residual moments. Analysis
continues only while every resulting coarse axis has at least five cells, and
synthesis restores the stored levels in reverse order. Induction on
\eqref{eq:fullcycle} then gives
\begin{equation}
 (\mathcal L^{(J)})^{-1}\mathcal L^{(J)}=I.
 \label{eq:multilevelinverse}
\end{equation}
For a one-dimensional fine line of length \(2N\), the \(N\) parent averages
and \(N\) residual moments are necessary and sufficient coordinates. For a
two-dimensional \(2N\times2M\) raster, the \(NM\) means and three \(NM\)
residual moment fields likewise account for every scalar degree of freedom.
% END CONSERVATIVE MULTIRESOLUTION CORE

\section{Direct Evaluation and Engineering Properties}

\subsection{Line analysis}

For each component of a source line, direct evaluation involves the following
fixed sequence.

\begin{enumerate}
\item Compute \(m_i,q_i\) from \eqref{eq:mint}--\eqref{eq:qbnd}.
\item Form \(P_{ij}\) and the raw current \(a_{ij}\).
\item Fill zero secant signs and locate the ordered boundaries
      \eqref{eq:candidates}--\eqref{eq:boundary}, using the exact cost
      recurrence \eqref{eq:ledgerrecurrence}.
\item Sort the five current breakpoints, locate the scalar interval, and
      apply \eqref{eq:thresholdcurrent}--\eqref{eq:thresholdroot}.
\item Integrate \(c_i\) once and evaluate each requested point by
      \eqref{eq:synthesis} or de Casteljau recursion.
\end{enumerate}

The boundary action compares at most nine candidate locations over at most
ten current slots. Equation~\eqref{eq:ledgerrecurrence} reduces the ordinary
cost scan from repeated ten-slot sums to at most ten initial slot evaluations
and eight signed cost updates, while the projection dimension remains five. A line with
\(N\) samples and \(C\) components therefore has \(O(NC)\) analysis cost,
\(O(NC)\) state, and constant cost per target evaluation. Consequently, the
compact form contains neither an image-sized matrix nor a global
eigendecomposition or convergence loop.

For comparison, the unprojected two-jet uses the same derivative stencils and
six quintic controls, but omits the ledger and scalar-threshold admission. A
direct Lanczos-3 evaluation uses at most six taps per Cartesian factor,
whereas EASU uses twelve two-dimensional fetches per target together with its
direction and anisotropic-weight calculations. Separable PCHIP first
constructs coupled nodal derivatives and then evaluates one cubic per factor.
CONV's additional work is therefore concentrated in the ordered-boundary
scan, one fixed five-key sort, at most six scalar interval tests, and five
clips per cell and component. Cartesian implementations of CONV, the raw
quintic, Lanczos, and PCHIP require a factor barrier; EASU instead evaluates a
direct two-dimensional target rule. Section~XI reports reference-CPU wall
times in order to expose the scale of this implementation cost.

\paragraph{Measured native cost reduction}
A matched native C implementation on an Apple M4 Mini measures the recurrence
against the original repeated-summation ledger, including allocations, basin
analysis, nodal geometry, and both Cartesian synthesis orders. With 31 warmed
repetitions and alternating implementation order, the four-worker median for
a cameraman \(512\to64\to512\) round trip decreases from 6.825 to
5.696\,ms, a 16.54\% time reduction. Across ten image, scale, and noise
workloads, the measured reductions are 10.63--16.77\% with four workers
and 13.30--20.35\% with one worker. Compilation is excluded. These are
complete-pipeline measurements for this implementation and host, not a
universal speed guarantee. Profile and image comparisons, including
quantized ties and wide dynamic ranges, agree; the one-worker audit also
checks the raw binary32 output bits. The mathematical reconstruction and
its approximation properties are unchanged.

\subsection{Parallel structure}

During two-jet formation, ledger construction, projection, and evaluation,
all lines of one Cartesian factor are independent and may be processed in
parallel. A barrier is required only between factors, since the next factor
analyzes the values produced by the preceding one. When several target
lattices query the same source, the arrays \(\{Q_{ij}\}\) or
\(\{c_{ij}\}\) may be cached.

\subsection{Boundary evaluation}

Equations \eqref{eq:mbnd} and \eqref{eq:qbnd} give the complete boundary
rule, so padded samples are not required. At a nested source node, the stored
observation may be returned directly. At a cell end, assigning the final
control as \(Q_{i5}=y_{i+1}\), while forming the interior controls from
cumulative admitted current, preserves cardinality independently of the
summation order used for those interior controls.

\subsection{Exact arithmetic and finite precision}

For \eqref{eq:ledgerrecurrence}, floating-point reassociation can affect
nearly tied costs. The native binary32 implementation accumulates squared
currents and recurrence updates in binary64, retaining the smallest and
second-smallest candidate costs, \(\widehat C_1\leq\widehat C_2\), and
\(\mathcal E_\ell=\sum_{k\in\mathcal N_\ell}a_k^2\). With
\(\varepsilon_{64}=2^{-52}\), it recomputes the candidate costs in their
original summation order whenever
\begin{equation}
 \widehat C_2-\widehat C_1
 \leq64\varepsilon_{64}\mathcal E_\ell,
 \label{eq:ledgertieguard}
\end{equation}
or the accumulated energy is nonfinite. The original smallest-index tie
rule is then applied. The bound covers the comparison uncertainty from
at most ten initial terms and eight signed updates in this binary32-to-binary64
realization. This is a numerical ambiguity safeguard for the existing
objective; it neither changes the admissible boundaries nor introduces a
geometric tolerance. The empirical bit comparisons above supplement this
local arithmetic safeguard and do not assert bitwise reversibility of the
complete multiresolution representation.


The identities in Sections VI--IX are algebraic. Direct floating-point
evaluation, however, incurs ordinary rounding error, and an approximate
numerical cycle does not prove \eqref{eq:fullcycle}. Bitwise reversible coding
therefore requires a reversible number representation: for integer data, the
lifting residual may be stored in an integer lifting realization; for
floating data, an error-free sum/difference expansion may be stored with the
detail when bit identity is required. Without that additional representation,
the theorem remains an exact-arithmetic statement, and ordinary
floating-point addition and subtraction are not claimed to be bitwise
identical.

% BEGIN COMPOSED SYNTHETIC EVIDENCE
\raggedbottom
\section{Synthetic Evaluation}

\subsection{Declared analytic design}

The evaluation uses a deterministic census of closed-form fields, without a
natural image, interpolated surrogate, trained parameter, or selected
favorable case. For a target field $f$, the source nodes form the exact
even--even sublattice, and error is evaluated on the remaining target nodes
after a four-target-node boundary crop:
\begin{equation}
 \operatorname{MSE}_{H}(\widehat f,f)
 =\frac{1}{|H|}\sum_{x\in H}|\widehat f(x)-f(x)|^2.
 \label{eq:empiricalmse}
\end{equation}
Thus the reference in \eqref{eq:empiricalmse} is the analytic field itself
rather than another reconstruction method.

The one-dimensional census contains 56 profiles at source lengths
$9,17,33,65,129$: 32 sinusoidal carriers from four frequencies and eight
phases, eight translated smooth transitions, eight Gaussian bumps, and eight
quadratic chirps. The two-dimensional census contains 28 fields at source
sides $9,17,33$: smooth transitions and windowed carriers at six off-axis
orientations and two subpixel phases, followed by two radial transitions and
two crossing-wave phases. The discontinuity census contains a step and a box
at 16 phases for each of five source lengths, giving 160 topology cases.

The one-dimensional census evaluates formal CONV and its raw two-jet
quintic before \(\Pi_{\mathcal C}\). The two-dimensional census evaluates
the bounded-work CONV$^{*}$ realization of Section~XII and audits it against
formal Eikonal CONV separately. Comparators are normalized endpoint-aligned
Lanczos-3, endpoint-aligned bilinear interpolation, and the EASU construction
specified below \cite{lanczos1956,lottes2021}. Every method uses its stated
fixed rule. Since the declared cases span many orders of magnitude, MSE is
aggregated by geometric mean, while paired win--loss counts retain the
individual case outcomes.

\subsection{Generator specification}

All analytic coordinates lie in \([0,1]\). For \(k=0,\ldots,7\), let
\(\phi_k=k\pi/4\). The one-dimensional families are
\begin{align}
 f_{\rm car}(x)&=0.5+0.38\sin(2\pi\nu x+\phi_k),\nonumber\\
 f_{\rm edge}(x)&=0.5+0.45\tanh((x-c)/w),\nonumber\\
 f_{\rm bump}(x)&=0.08+0.84\exp(-((x-c)/\sigma)^2),\nonumber\\
 f_{\rm chirp}(x)&=0.5+0.34\sin(2\pi(0.5x+2.5x^2)+\phi_k).
 \label{eq:onedgenerators}
\end{align}
Here \(\nu\in\{0.75,1.5,2.25,3\}\),
\(c\in\{0.44,0.48,0.52,0.56\}\),
\(w\in\{0.025,0.05\}\), and
\(\sigma\in\{0.045,0.09\}\). The carrier uses all \((\nu,\phi_k)\)
pairs; edge and bump use all stated \((c,w)\) and \((c,\sigma)\) pairs;
the chirp uses all \(\phi_k\).

For the two-dimensional fields, put
\begin{align}
 p_\theta&=x\cos\theta+y\sin\theta,\nonumber\\
 n_\theta&=(x-\tfrac12)\cos\theta+(y-\tfrac12)\sin\theta,\nonumber\\
 e(x,y)&=\sin^2(\pi x)\sin^2(\pi y),\nonumber\\
 r&=\sqrt{(x-0.47)^2+(y-0.53)^2}.
\end{align}
With \(h=(N-1)^{-1}\), the four families are
\begin{align}
 F_{\rm edge}&=0.5+0.43\tanh((n_\theta-\delta)/0.035),\nonumber\\
 F_{\rm car}&=0.5+0.34e\sin(2\pi(3.25)p_\theta+\phi),\nonumber\\
 F_{\rm curved}&=0.5+0.42\tanh((r-0.27-\delta)/0.032),\nonumber\\
 F_{\rm cross}&=0.5+0.19\{g_1+g_2\},\nonumber\\
 g_1&=\sin(2\pi(2.75x+0.9y)+\phi),\nonumber\\
 g_2&=\sin(2\pi(-0.9x+2.75y)-\phi).
 \label{eq:twodgenerators}
\end{align}
Here
\(\theta\in\{7.5,22.5,37.5,52.5,67.5,82.5\}^{\circ}\). Edge fields
use \(\delta=\pm h/4\), carriers use \(\phi\in\{0,\pi/2\}\), curved
fields use \(\delta=\pm h/4\), and crossing fields use
\(\phi\in\{0,\pi/2\}\).

For a source line of length \(N\), topology phase
\(s_k=(k+1/2)/16\) gives
\begin{align}
 y_j^{\rm step}&=\mathbf 1\{j\geq N/2-1/2+s_k\},\nonumber\\
 y_j^{\rm box}&=\mathbf 1\{|j-(N/2-1/2+s_k)|\leq N/6\}.
 \label{eq:topologygenerators}
\end{align}
Here \(k=0,\ldots,15\).
Equations \eqref{eq:onedgenerators}--\eqref{eq:topologygenerators}, the
listed grids, and the held-out rule in \eqref{eq:empiricalmse} determine the
entire population without an external corpus.

\subsection{Comparator definitions}

Let \(\operatorname{sinc}(t)=\sin(\pi t)/(\pi t)\) and
\(L_3(t)=\operatorname{sinc}(t)\operatorname{sinc}(t/3)\) for
\(|t|<3\), zero otherwise. Resampling \(n\) nodes to \(m\) nodes uses
\begin{align}
 x_r&=\frac{r(n-1)}{m-1},\nonumber\\
 \rho&=\min\!\left\{1,\frac{m-1}{n-1}\right\},\nonumber\\
 w_{rj}&=\frac{L_3(\rho(x_r-j))}{\sum_q L_3(\rho(x_r-q))}.
 \label{eq:lanczosanalysis}
\end{align}
Only in-domain \(j\) are summed; hence the boundary rule is support
truncation followed by renormalization. This gives the endpoint-aligned
Lanczos arm for the point-sample interpolation census. Bilinear uses the same
endpoint coordinates with support restricted to the two neighboring nodes,
whereas the cell-centered conservative transform ablation is stated
separately in
\eqref{eq:celllanczoscoordinate}--\eqref{eq:celllanczosweight}.

The EASU arm is a binary64 CPU transcription of the FSR 1.0 EASU
construction. Scalar fields are replicated to RGB. It uses the twelve EASU
fetch locations, exact binary64 reciprocal and square root rather than the
shader approximations, and no RCAS pass. Target \((r,s)\) is evaluated at
source coordinate \((r/2,s/2)\); this is the complete coordinate-constant
specialization for the nested endpoint-aligned grid. Out-of-domain fetches
are clamped to the nearest source pixel. The raw two-jet arm uses the same
stencils, quintic controls, Cartesian order, and evaluation rule as CONV but
sets \(c_i=a_i\), omitting the ledger and \(\Pi_{\mathcal C}\).
PCHIP appears only in the timing audit; SciPy's \texttt{PchipInterpolator}
is applied horizontally and then vertically on the same endpoint-aligned
nested coordinates, without padding.

\begin{figure*}[t]
 \centering
 \includegraphics[width=\textwidth]{conv_paper_assets/convergence.pdf}
 \caption{Resolution census against direct analytic target values. Each point
 is the geometric mean over the complete declared population at that source
 size. The coarse regime exposes underresolution; the resolved regime exposes
 the approximation order of each fixed operator.}
 \label{fig:convconvergence}
\end{figure*}

\subsection{Resolution and paired fidelity}

Table~\ref{tab:pairedfidelity} compares formal one-dimensional CONV and
two-dimensional CONV$^{*}$ with Lanczos-3, such that a ratio below one favors
the corresponding CONV form. At nine nodes neither population is resolved,
and Lanczos-3 has the lower aggregate error. From 17 nodes onward, however,
CONV wins the majority of one-dimensional cases; at 129 nodes, it wins all
56. In two dimensions, CONV$^{*}$ has the lower geometric-mean error at 17
and 33 nodes. The paired count at 33 nodes is even because Lanczos retains a
small advantage on steep transition fields, whereas CONV$^{*}$ has larger
advantages on the resolved carrier and crossing fields. The aggregate ratio
therefore records the magnitude of this distinction without discarding its
paired count.

\begin{table}[H]
\caption{Paired point-fidelity results versus Lanczos-3}
\label{tab:pairedfidelity}
\centering
\scriptsize
\setlength{\tabcolsep}{3.0pt}
\begin{tabular}{c|cc|cc}
\hline
Source & \multicolumn{2}{c|}{1-D point} &
\multicolumn{2}{c}{2-D point} \\
nodes & W--L & MSE ratio & W--L & MSE ratio \\
\hline
9   & 28--28 & 1.247 & 9--19  & 1.277 \\
17  & 36--20 & 0.0929 & 19--9 & 0.800 \\
33  & 44--12 & 0.00751 & 14--14 & 0.126 \\
65  & 52--4  & 0.000481 & -- & -- \\
129 & 56--0  & 0.0000275 & -- & -- \\
\hline
\end{tabular}
\end{table}

The EASU arm establishes a different engineering point. At the $9\times9$
source size it has the lowest two-dimensional geometric-mean MSE,
$2.38\times10^{-3}$. Its nested-source maximum error is $1.81\times10^{-1}$,
whereas CONV$^{*}$, Lanczos-3, and bilinear are cardinal to their stated
floating-point precision. At $33\times33$, CONV$^{*}$ gives
$4.31\times10^{-8}$ geometric-mean MSE, compared
with $3.42\times10^{-7}$ for Lanczos-3, $3.23\times10^{-5}$ for bilinear,
and $4.67\times10^{-5}$ for EASU. These measurements separate interpolation
fidelity from a reconstruction rule that is permitted to alter source nodes.

\subsection{Ordered-current ablation}

The raw two-jet arm isolates the admission map by retaining the compact
quintic proposal and removing only \(\Pi_{\mathcal C}\). From
Table~\ref{tab:projectionablation}, it can be seen that admission does not
obtain its topology result by replacing the high-order proposal: aggregate
MSE is slightly lower on every reported grid and becomes indistinguishable
from the raw proposal in the resolved regime.

\begin{table}[H]
\caption{Raw two-jet versus admitted CONV geometric-mean MSE; the 2-D columns use CONV$^{*}$}
\label{tab:projectionablation}
\centering
\scriptsize
\setlength{\tabcolsep}{4.0pt}
\begin{tabular}{c|cc|cc}
\hline
Source & \multicolumn{2}{c|}{1-D} & \multicolumn{2}{c}{2-D} \\
nodes & Raw & CONV & Raw & CONV \\
\hline
9  & $2.95\!\times\!10^{-4}$ & $2.50\!\times\!10^{-4}$ &
     $4.02\!\times\!10^{-3}$ & $3.87\!\times\!10^{-3}$ \\
17 & $7.02\!\times\!10^{-7}$ & $6.88\!\times\!10^{-7}$ &
     $2.63\!\times\!10^{-5}$ & $2.48\!\times\!10^{-5}$ \\
33 & $5.51\!\times\!10^{-10}$ & $5.49\!\times\!10^{-10}$ &
     $4.31\!\times\!10^{-8}$ & $4.31\!\times\!10^{-8}$ \\
\hline
\end{tabular}
\end{table}

The distinction is decisive on the discontinuity census. The raw quintic
has \(S>0\) in all 160 step/box cases and reaches a surplus of eight sign
changes; CONV has \(S=0\) in all 160. Thus the compact jet supplies the
resolved-field accuracy, while ordered-current admission removes the
unwitnessed variation and, on this census, does not increase aggregate MSE.

\subsection{Topology and excursion}

For the discontinuity census, the measured topology quantity is
\begin{equation}
 S=\max\{0,\var(d_h\widehat f)-\var(d_cf)\}.
 \label{eq:measuredsurplus}
\end{equation}
CONV has $S=0$ in all 160 cases, agreeing with
\eqref{eq:vardim}; bilinear also has zero surplus. The raw quintic and
Lanczos-3 each have $S>0$ in all 160 cases and reach eight additional sign
changes. This isolates the effect of admission from that of the compact jet
and distinguishes CONV from a sharp signed-lobe kernel.
CONV's maximum amplitude excursion on these step and box cases is 0.1008,
compared with 0.1194 for Lanczos-3 and zero for bilinear. The zero-surplus
result therefore does not conceal the separately measured motion of an
existing extremum.

\begin{figure*}[t]
 \centering
 \includegraphics[width=0.58\textwidth]{conv_paper_assets/topology.pdf}
 \caption{One-dimensional discontinuity topology. The plot counts every
 declared step or box phase that creates a derivative-sign pair absent from
 the samples.}
 \label{fig:convtopology}
\end{figure*}

% BEGIN CONSERVATIVE MULTIRESOLUTION EVIDENCE
The conservative zero-detail transform ablation uses the dyadic state of
Section~IX. Fine
arrays have 16, 32, or 64 cells per axis. The one-dimensional census contains
one Nyquist sequence, eight phases of a $0.92\pi$ carrier, and eight phases
of a quadratic chirp at each size. The two-dimensional census contains axial
and checkerboard Nyquist arrays, eight half-plane interfaces over four angles
and two offsets, and four chirp phases. Half-plane pixels are their analytic
polygonal cell coverage; no supersampling constructs the reference.

Each method uses its own analysis and synthesis. CONV-C denotes
\eqref{eq:scaleaverage} followed by zero-residual synthesis
\eqref{eq:synthleft}--\eqref{eq:synthright} and its two-dimensional moment
form. Box-bilinear uses the same positive analysis and pixel-center linear
synthesis. For Lanczos-3, if \(\rho=M/N\), the target center is
\begin{equation}
 u_j=\frac{j+1/2}{\rho}-\frac12,\qquad
 \kappa=\min(1,\rho),
 \label{eq:celllanczoscoordinate}
\end{equation}
and its normalized in-domain weight is
\begin{equation}
 \begin{aligned}
 w_{jn}&\propto
 \operatorname{sinc}\!\left(\kappa(u_j-n)\right)
 \operatorname{sinc}\!\left(\frac{\kappa(u_j-n)}3\right),\\
 &\hspace{25mm}|\kappa(u_j-n)|<3.
 \end{aligned}
 \label{eq:celllanczosweight}
\end{equation}
The same definition is used in analysis and synthesis; finite boundaries are
truncated and renormalized. MSE is evaluated against the original fine cell
array after a four-cell boundary crop.

\begin{figure*}[t]
 \centering
 \includegraphics[width=\textwidth]{conv_paper_assets/conservative_cycle.pdf}
 \caption{Conservative zero-detail transform ablation. Left: geometric-mean MSE
 over all 14 two-dimensional cases at each size. Right: mean MSE over the
 eight polygonal diagonal interfaces. Each curve uses its method's declared
 analysis and synthesis.}
 \label{fig:conservativecycle}
\end{figure*}

\begin{table}[H]
\caption{Conservative two-dimensional zero-detail ablation}
\label{tab:conservativecycle}
\centering
\scriptsize
\setlength{\tabcolsep}{3.0pt}
\begin{tabular}{c|ccc|c}
\hline
Fine side & CONV-C & L3 & Box-bilin. & Edge ratio \\
\hline
16 & $2.78\!\times\!10^{-2}$ & $3.16\!\times\!10^{-2}$ &
$4.57\!\times\!10^{-2}$ & 0.845 \\
32 & $1.45\!\times\!10^{-2}$ & $1.64\!\times\!10^{-2}$ &
$2.35\!\times\!10^{-2}$ & 0.807 \\
64 & $8.66\!\times\!10^{-3}$ & $9.86\!\times\!10^{-3}$ &
$1.41\!\times\!10^{-2}$ & 0.793 \\
\hline
\end{tabular}
\end{table}

CONV-C has the lowest aggregate MSE at all three sizes. Its mean diagonal
interface MSE is 15.5, 19.3, and 20.7 percent below Lanczos-3 as size
increases. The chirp comparison is phase dependent: CONV-C has the lower
four-phase mean at side 16, while Lanczos-3 is lower by 0.50 and 1.14 percent
at sides 32 and 64. Thus the aggregate result does not conceal the remaining
high-frequency case.

The conservative restriction maps axial Nyquist and checkerboard inputs to
zero exactly, agreeing with \eqref{eq:nyquistresponse} and
\eqref{eq:nyquist2d}. CONV-C has zero horizontal and vertical sign surplus in
all 42 two-dimensional cases. Six retained-detail audits---one line and one
three-channel raster at each size---have maximum binary64 round-trip residual
$8.88\times10^{-16}$; the exact identity is proved in Section~IX.

The preceding census isolates the reversible conservative state. The
technical fields in Fig.~\ref{fig:currentmatchedplate} instead use the current
principal raster operator in a matched
\(64^2\!\rightarrow32^2\!\rightarrow64^2\) cycle. CONV$^{*}$ applies basin
restriction followed by its present \(Q_1\)-compiled synthesis. Lanczos-3
uses its endpoint-aligned, scale-widened kernel in both
directions, while bilinear uses endpoint-aligned linear interpolation in both
directions. Thus no row inherits the CONV-C zero-detail synthesis shown in
Fig.~\ref{fig:conservativecycle}.

\begin{figure*}[t]
 \centering
 \includegraphics[width=0.92\textwidth]{conv_paper_assets/current_matched_plate.pdf}
 \caption{Fine-grid truth and current matched resize cycles. For the selected
 diagonal interface, MSE is $1.86\times10^{-3}$ for CONV$^{*}$,
 $1.37\times10^{-3}$ for Lanczos-3, and $2.64\times10^{-3}$ for bilinear;
 the corresponding range excursions are $3.91\times10^{-2}$,
 $9.63\times10^{-2}$, and zero. On the quadratic chirp, MSE is
 $5.37\times10^{-2}$, $5.20\times10^{-2}$, and $6.00\times10^{-2}$ in the
 same order.}
 \label{fig:currentmatchedplate}
\end{figure*}
% END CONSERVATIVE MULTIRESOLUTION EVIDENCE

\subsection{Technical synthetic fields}

Figure~\ref{fig:convtechnical} contains one declared case from each
two-dimensional family at source side 17. The images use a dense
$513\times513$ endpoint-aligned query lattice; its comparison grid retains
the scored target side 33 and reports held-out MSE, maximum excursion, and
nested-source error.
Lanczos is slightly better on the selected edge and radial transition;
CONV$^{*}$ is better on the carrier and crossing field and has the smaller maximum
excursion while remaining cardinal. The tabulated raw row records the
projection's local effect on the same fields.

\begin{figure*}[t]
 \centering
 \includegraphics[width=0.90\textwidth]{conv_paper_assets/technical_synthetics.pdf}
 \vspace{1mm}

 \scriptsize
 \setlength{\tabcolsep}{4.0pt}
 \begin{tabular}{l|cccc|cc}
 \hline
 Method & Edge MSE & Carrier MSE & Curved MSE & Crossing MSE &
 Max. excursion & Source $L_\infty$ \\
 \hline
 CONV$^{*}$ & $1.09\!\times\!10^{-4}$ & $1.43\!\times\!10^{-6}$ & $7.93\!\times\!10^{-4}$ &
 $9.98\!\times\!10^{-7}$ & $1.92\!\times\!10^{-2}$ & $2.98\!\times\!10^{-8}$ \\
 Raw 2-jet & $1.22\!\times\!10^{-4}$ & $1.43\!\times\!10^{-6}$ & $8.68\!\times\!10^{-4}$ &
 $9.98\!\times\!10^{-7}$ & $2.49\!\times\!10^{-2}$ & 0 \\
 Lanczos-3 & $9.46\!\times\!10^{-5}$ & $4.56\!\times\!10^{-6}$ & $7.48\!\times\!10^{-4}$ &
 $7.89\!\times\!10^{-6}$ & $3.31\!\times\!10^{-2}$ & $1.11\!\times\!10^{-16}$ \\
 Bilinear & $5.96\!\times\!10^{-4}$ & $3.29\!\times\!10^{-4}$ & $1.97\!\times\!10^{-3}$ &
 $5.34\!\times\!10^{-4}$ & 0 & 0 \\
 EASU & $4.26\!\times\!10^{-4}$ & $1.38\!\times\!10^{-4}$ & $7.34\!\times\!10^{-4}$ &
 $2.31\!\times\!10^{-4}$ & 0 & $9.05\!\times\!10^{-2}$ \\
 \hline
 \end{tabular}
 \caption{Technical synthetic fields and their local comparison grid. Images
 use a dense endpoint-aligned query lattice and one common $[0,1]$ display
 scale; the table reports unquantized binary64 values at target side 33
 against direct analytic evaluation.}
 \label{fig:convtechnical}
\end{figure*}

\subsection{Eikonal compilation conformance}

CONV$^{*}$ evaluates the two formal Cartesian factor orders and replaces the
all-source Eikonal coefficient \(\beta\) by the bilinear extension
\(\beta^*=I_h\beta\) of its cardinal nodal trace. Subtraction from the formal
operator gives the exact relation
\begin{equation}
 \Tcal^*Y-\Tcal Y
 = (\beta^*-\beta)(\Tcal_{yx}Y-\Tcal_{xy}Y).
 \label{eq:evidenceimplementationdefect}
\end{equation}
Therefore the fast form is exact wherever the factor orders commute; away
from that set, its complete discrepancy is precisely the product in
\eqref{eq:evidenceimplementationdefect}.

Table~\ref{tab:eikonalconformance} evaluates that discrepancy on 16 fixed
edge, carrier, curved-interface, and crossing fields. Every row has the same
$65\times65$ target lattice; only source density and scale change. RMS is
formed over all target values and cases. The maximum column is a census
maximum, not a confidence estimate. Across the three rows, replacing the
formal action changes geometric-mean truth MSE by less than 0.01 percent.

\begin{table}[H]
\caption{CONV$^{*}$ conformance to formal Eikonal CONV at target side 65}
\label{tab:eikonalconformance}
\centering
\scriptsize
\setlength{\tabcolsep}{2.5pt}
\resizebox{\columnwidth}{!}{%
\begin{tabular}{c|c|ccc|cc}
\hline
Source & Scale & RMS & Median $L_\infty$ & Max. $L_\infty$ &
 Formal ms & CONV$^{*}$ ms \\
\hline
9  & 8 & $1.75\!\times\!10^{-4}$ & $1.50\!\times\!10^{-3}$ &
 $4.63\!\times\!10^{-3}$ & 11.0 & 0.402 \\
17 & 4 & $5.67\!\times\!10^{-5}$ & $6.05\!\times\!10^{-4}$ &
 $2.43\!\times\!10^{-3}$ & 43.6 & 0.460 \\
33 & 2 & $5.95\!\times\!10^{-6}$ & $1.56\!\times\!10^{-5}$ &
 $8.50\!\times\!10^{-4}$ & 236 & 0.659 \\
\hline
\end{tabular}
}
\end{table}

\subsection{Aggregate Green admission audit}

The optional Green coordinate \eqref{eq:greenadmission} was evaluated on the
same 16 fixed fields using the source-derived determinant-one metric,
\(\rho=1\), no-flux boundaries, a monotone anisotropic Selling stencil, and
the cardinal weighted-harmonic form \eqref{eq:greenharmonic}
\cite{mirebeau2014}. Table
\ref{tab:greenadmission} compares the result with the local coordinate
\(\beta^*=I_h\beta\). Although the coordinate changes materially, the exact
commutator identity \eqref{eq:greenblenddefect} suppresses most of that
difference in the reconstructed field.

\begin{table}[H]
\caption{Cardinal Green admission versus local \(Q_1\) compilation}
\label{tab:greenadmission}
\centering
\scriptsize
\setlength{\tabcolsep}{2.2pt}
\resizebox{\columnwidth}{!}{%
\begin{tabular}{c|ccc|cc}
\hline
Source & Local MSE & Green MSE & Ratio & \(\beta\) RMS & Output RMS \\
\hline
\(9^2\)  & \(2.3888571\!\times\!10^{-3}\) &
 \(2.3890219\!\times\!10^{-3}\) & 1.000069 &
 \(4.76\!\times\!10^{-2}\) & \(7.78\!\times\!10^{-5}\) \\
\(17^2\) & \(1.8625653\!\times\!10^{-4}\) &
 \(1.8626100\!\times\!10^{-4}\) & 1.000024 &
 \(3.89\!\times\!10^{-2}\) & \(8.50\!\times\!10^{-6}\) \\
\hline
\end{tabular}%
}
\end{table}

Because its discrete diagonal is finite, the raw finite-grid ratio \(V/U\)
has maximum source error 0.542 and 0.628 on the two lattices. Imposing the
nodal traces makes that error exactly zero. Consequently, the global
coordinate supplies a closed metric-harmonic option and a formal comparator;
however, this census yields no fidelity reason to displace the local \(Q_1\)
compilation.

\subsection{Native CPU cost}

Table~\ref{tab:nativecpu} reports three-channel float32 wall time on an M4
Mac mini. Native CONV$^{*}$, Lanczos-3, bilinear, and EASU use ARM
NEON and the same ten-worker pool; PCHIP is SciPy's compiled implementation.
Each entry is the median of 21 calls after three warmups, without affinity.
The largest median absolute deviation in the table is 0.065 ms. Fidelity and
timing use the same formulas, but the EASU timing arm uses the native shader
translation rather than the binary64 audit port.

\begin{table}[H]
\caption{Native RGB point-synthesis time to $512^2$ (ms)}
\label{tab:nativecpu}
\centering
\scriptsize
\setlength{\tabcolsep}{2.7pt}
\begin{tabular}{l|rr}
\hline
Method & from $64^2$ & from $256^2$ \\
\hline
CONV$^{*}$ & 2.985 & 9.725 \\
Lanczos-3  & 2.045 & 3.147 \\
Bilinear   & 0.282 & 0.834 \\
EASU       & 1.715 & 1.704 \\
PCHIP      & 3.750 & 15.131 \\
\hline
\end{tabular}
\end{table}

The native result places the cost rather than hiding it: at twofold synthesis
CONV$^{*}$ is 3.09 times the Lanczos-3 time and 5.71 times the EASU time. At
eightfold synthesis it is 1.46 times the Lanczos-3 time and 1.74 times the
EASU time. These measurements
characterize the present CPU realization;
they are not used in any analytical claim.

Table~\ref{tab:binary64audit} is an implementation audit, not a substitute for
the preceding exact-arithmetic proofs. It uses 80 deterministic Gaussian
three-channel lines over five lengths. The current-mass, cardinal, and lifting
residuals remain at binary64 rounding scale, and no admitted current has sign
variation exceeding its observed secant sequence.

\begin{table}[H]
\caption{Binary64 conformance audit over 80 deterministic trials}
\label{tab:binary64audit}
\centering
\scriptsize
\begin{tabular}{l|c}
\hline
Quantity & Maximum observed value \\
\hline
Cell-current mass residual & $8.88\times10^{-16}$ \\
Nested cardinal error & $1.33\times10^{-15}$ \\
Retained-detail lifting error & $4.44\times10^{-16}$ \\
Sign-variation surplus & 0 \\
\hline
\end{tabular}
\end{table}

The executable generator, complete per-case records, and figure construction
form the supplementary material accompanying the paper. Aggregate curves and
tables are generated from those records without case removal or post hoc
weighting; the analytic population itself is fully specified above.
% END COMPOSED SYNTHETIC EVIDENCE

% BEGIN COMPOSED CONVSTAR
\section{Can CONV Be Made Efficient Using Convolution?}
\label{sec:convstar}

Because CONV specifies the two-jet, lineage ledger, signed-fibre projection,
Bernstein synthesis, and Eikonal factor blend independently of an execution
backend, the same operator admits a bounded-work FIR realization. Denoted by
CONV$^{*}$, this realization eliminates the one-dimensional jet algebra into
fixed convolution banks, batches the exact projection, compiles repeated
phases into a polyphase output bank, and replaces the all-source action sum by
the local nodal compilation derived below. The asterisk therefore denotes a
computational refinement of the same admitted fibre, rather than another
interpolation model. In two dimensions, the FIR bank replaces each
one-dimensional occurrence of \(\mathcal C\) in \(\Tcal_{xy}\) and
\(\Tcal_{yx}\); at a fixed endpoint-aligned scale, the same polyphase bank is
used in opposite factor orders and the resulting proposals are blended
pointwise, without forming a source-specific frame or chart bank.

\subsection{Local compilation of Eikonal admission}

Let \(\eta_p\) be the nodal order coordinate in
\eqref{eq:witnesscompatibility}. A local compilation is obtained by replacing
the all-source action interpolant \(\beta\) with its bilinear nodal extension,
\begin{align}
 \beta^*(x)&=\sum_{p\in\operatorname{cell}(x)}b_p(x)\eta_p,\nonumber\\
 \Tcal^*Y&=(1-\beta^*)\Tcal_{xy}Y+\beta^*\Tcal_{yx}Y.
 \label{eq:starlocaladmission}
\end{align}
\begin{proposition}[Cardinal compilation identity]
Let \(I_h\) denote cellwise bilinear \(Q_1\) interpolation on the source
lattice. Then
\begin{equation}
 \boxed{\beta^*=I_h\beta.}
 \label{eq:starcardinalcompilation}
\end{equation}
Consequently, \(\beta^*\) is the unique continuous cellwise bilinear extension
of the formal Eikonal admission's nodal trace.
\end{proposition}

\begin{proof}
The Kronecker action identity \eqref{eq:kroneckeraction} gives
\(\beta(x_p)=\eta_p\) at every source node. Applying \(I_h\) therefore yields
\((I_h\beta)(x)=\sum_{p\in\operatorname{cell}(x)}b_p(x)\beta(x_p)\), which is
exactly \eqref{eq:starlocaladmission}.
\end{proof}

Thus the compilation is not an independently chosen approximation rule, but
the canonical \(Q_1\) interpolant of the formal admission data; it is
parameter-free, bounded in \([0,1]\), and permutation covariant. Subtracting
\eqref{eq:witnessedblend} from \eqref{eq:starlocaladmission} gives the exact
implementation defect
\begin{equation}
 \Tcal^*Y-\Tcal Y
 = (\beta^*-\beta)(\Tcal_{yx}Y-\Tcal_{xy}Y),
 \label{eq:stardefect}
\end{equation}
and therefore, componentwise,
\begin{equation}
 |\Tcal^*Y-\Tcal Y|
 \leq |\beta^*-\beta|\,|\Tcal_{yx}Y-\Tcal_{xy}Y|.
 \label{eq:stardefectbound}
\end{equation}
Therefore the local compilation is exact wherever the factor orders commute,
including constant and affine coordinate fields. Moreover,
\eqref{eq:stardefect} isolates the sole two-dimensional approximation
introduced by CONV$^{*}$, since the current construction within each factor
remains unchanged.

For direct evaluation, \(\eta_p\) requires only the squared first-jet
channels already produced by the two Cartesian front ends:
\begin{equation}
 \eta_p=
 \frac{\sum_c(D_yY_p^{(c)})^2}
 {\sum_c\bigl[(D_xY_p^{(c)})^2+(D_yY_p^{(c)})^2\bigr]},
 \label{eq:stardirecteta}
\end{equation}
with the value one half when the denominator vanishes. Consequently, neither
an eigendecomposition nor the cross term of the current tensor is required
to evaluate the order coordinate. At a target point, the four nodal values
in \eqref{eq:starlocaladmission} may be combined with the final convex blend
in the same output write. This evaluates \(I_h\eta\) exactly while avoiding a
separate target-sized coordinate raster and its two Cartesian interpolation
passes.

\subsection{Aggregate Green admission}

A global elliptic continuation of the same nodal coordinate is obtained as
follows. Put \(K=M^{-1}\), let \(\rho\) be the source density per unit
Riemannian area, and define
\begin{equation}
 \begin{aligned}
  (\rho-\Delta_M)U&=\sum_a\delta_{s_a},\\
  (\rho-\Delta_M)V&=\sum_a\eta_a\delta_{s_a},\\
  \beta_G&=V/U.
 \end{aligned}
 \label{eq:greenadmission}
\end{equation}
subject to the no-flux boundary condition
\(n^TK\nabla U=n^TK\nabla V=0\). In the unit source-lattice coordinates used
throughout the paper, \(\rho=1\). Under \(x\mapsto cx\), both \(\rho\) and
\(\Delta_M\) scale by \(c^{-2}\), so the screening length is fixed by the
sampling density rather than by a fitted bandwidth. Moreover, the two
right-hand sides use one SPD operator and one factorization, without forming
an individual distance field for each source.

\begin{proposition}[Green transport coordinate]
The continuum coordinate in \eqref{eq:greenadmission} is cardinal by
continuous extension, satisfies
\begin{equation}
 \min_a\eta_a\leq\beta_G(x)\leq\max_a\eta_a,
 \label{eq:greenmaximum}
\end{equation}
and, away from the source sites, obeys
\begin{equation}
 \operatorname{div}\!\left(U^2M^{-1}\nabla\beta_G\right)=0.
 \label{eq:greenharmonic}
\end{equation}
Equivalently, it minimizes
\begin{equation}
 \mathcal E_U[\vartheta]
 =\frac12\int_\Omega U^2\nabla\vartheta^TM^{-1}\nabla\vartheta\,
 d\operatorname{vol}_M
 \label{eq:greenenergy}
\end{equation}
subject to the nodal traces \(\vartheta(s_a)=\eta_a\).
\end{proposition}

\begin{proof}
For \(\rho>0\), the shifted Neumann form is coercive and its Green kernel is
strictly positive. Consequently, \(V/U\) is a positive normalized
Green-kernel blend of the \(\eta_a\), which proves
\eqref{eq:greenmaximum} \cite{gilbargtrudinger2001}. In two dimensions, the
diagonal singularity is
\begin{equation}
 G_\rho(x,s_a)=-\frac{1}{2\pi}\log d_M(x,s_a)+O(1).
 \label{eq:greensingu}
\end{equation}
and therefore the self term dominates at \(s_a\), giving
\(\lim_{x\to s_a}\beta_G(x)=\eta_a\). Away from the sources, substituting
\(V=U\beta_G\) into the two homogeneous resolvent equations cancels the
common mass term; the product rule then yields \eqref{eq:greenharmonic}, whose
weak form is the Euler equation of \eqref{eq:greenenergy}.
\end{proof}

This construction continues the same nodal order coordinate, but it does not
fuse the inverse-square Eikonal admission exactly. Indeed,
\eqref{eq:greensingu} is logarithmic, whereas
\eqref{eq:actionpartition} uses \(d_M^{-2}\); hence no second-order local
elliptic Green function in two dimensions can make the two kernels
identical. Since a finite-stencil diagonal also removes the continuum
singularity, the raw discrete ratio \(V/U\) is not exactly cardinal. Exact
discrete cardinality is obtained by solving \eqref{eq:greenharmonic} with the
nodal traces imposed.

Both the global Green admission and the local compilation therefore have the
trace \(\eta_p\). The latter is the unique cellwise \(Q_1\) continuation
satisfying locality, linearity in the four nodal values, cardinality, and
separate affine exactness, whereas the former is the global minimizer of
\eqref{eq:greenenergy}. Their difference in the reconstructed field is
isolated exactly by
\begin{equation}
 \Tcal_{\beta_G}Y-\Tcal_{\beta^*}Y
 =(\beta_G-\beta^*)(\Tcal_{yx}Y-\Tcal_{xy}Y).
 \label{eq:greenblenddefect}
\end{equation}
Consequently, Green admission supplies a global metric-harmonic coordinate
when that property is required, while \(Q_1\) remains the compact
CONV$^{*}$ realization.

One Cartesian factor has the dataflow
\begin{equation}
 y\xrightarrow{\rm FIR}(a,\delta)
 \xrightarrow{\rm scan}\sigma
 \xrightarrow{\rm tensorized\ projection}c
 \xrightarrow{\rm polyphase}\widehat y.
 \label{eq:starflow}
\end{equation}
Since the factors retain CONV's horizontal--vertical barrier, the second
factor analyzes the values synthesized by the first. No single separable
two-dimensional convolution therefore represents the complete nonlinear
composition.

\subsection{Compact FIR current bank}

On an interior cell, let \(v_i=(y_{i+r})_{r=-2}^{3}\). Eliminating the
intermediate jets and Bezier controls yields
\begin{equation}
 a_i=Kv_i,\quad K=\frac1{240}
 \left[\begin{smallmatrix}
 4&-32&0&32&-4&0\\
 3&-16&-30&48&-5&0\\
 -7&39&-130&130&-39&7\\
 0&5&-48&30&16&-3\\
 0&4&-32&0&32&-4
 \end{smallmatrix}\right].
 \label{eq:starfir}
\end{equation}
Hence the complete interior proposal is a fixed bank of five six-tap FIR
filters. An equivalent realization first computes the two five-tap jet
channels and the secant, and then forms the five currents by pointwise
algebra. The latter requires fewer arithmetic operations on a conventional
CPU, whereas the fused form maps directly to a multi-output convolution
primitive. At the outer and near-outer cells, the stated second-order closure
is retained rather than inferred from \eqref{eq:starfir} or supplied by zero
padding.

\subsection{The irreducible sign scan}

Starting from the sampled signs
\(s_i=\operatorname{sign}(\delta_i)\), CONV$^{*}$ fills a zero from the last
witnessed nonzero sign, using the first future witness only for an initial
zero run. Applying the same local least-disagreement rule to the resulting
transition sites gives a segmented prefix scan followed by a short causal
transition scan.

\begin{proposition}[No bounded convolutional ledger]
No fixed finite-receptive-field operator can reproduce the zero-sign filling
rule for every line length.
\end{proposition}

\begin{proof}
Fix a receptive-field radius \(R\) and a zero secant at a site whose previous
\(R\) secants are also zero. Construct two lines with identical secants on
that complete receptive field, but place a positive witness immediately
before it in one line and a negative witness there in the other. The formal
ledger assigns opposite signs at the selected zero site. A bounded local
operator receives identical inputs and must return the same value, a
contradiction.
\end{proof}

Therefore the scan must be retained rather than approximated by a bounded
convolution. It involves linear work with one sign of persistent state per
component, while all lines of a Cartesian factor remain independent.

\subsection{Vectorized scalar-threshold admission}

CONV$^{*}$ batches the exact scalar projection
\eqref{eq:thresholdcurrent}--\eqref{eq:thresholdroot} by using a fixed sorting
network for the five breakpoints in each cell. The six resulting affine
intervals are tested in parallel, after which five signed clips return the
admitted current. Chunking consequently bounds temporary storage
independently of image size, and neither a combinatorial search nor an
iterative numerical solve is involved.

\subsection{Polyphase synthesis}

Let \(B_j^5(u)\) denote the degree-five Bernstein basis. Eliminating the
integrated controls gives
\begin{equation}
 \widehat y(i+u)=y_i+\sum_{k=0}^{4}g_k(u)c_{ik},\qquad
 g_k(u)=\sum_{j=k+1}^{5}B_j^5(u).
 \label{eq:starpolyphase}
\end{equation}
For integer refinement \(r\), only the phases \(u=p/r\),
\(p=1,\ldots,r-1\), are required. Their \(g_k\) values form a fixed
polyphase matrix. At twofold refinement,
\begin{equation}
 \widehat y(i+\tfrac12)=y_i+\frac{1}{32}
 [31,26,16,6,1]c_i.
 \label{eq:starmidpoint}
\end{equation}
Thus neither the Bezier controls nor their cumulative current sums need be
stored. Nested source nodes are copied directly for point interpolation,
whereas conservative multiresolution compiles the half-cell functional
\eqref{eq:halfcells}, block averages, and centered moment channels of
Section~IX; it does not replace them by strided anchor selection.
Accordingly, \eqref{eq:starmidpoint} is a compiled twofold phase rather than
the domain of CONV. For an arbitrary target length \(M\), the continuous
operator evaluates \eqref{eq:starpolyphase} at the coordinates
\eqref{eq:targetmap}; whenever the target uses a repeated finite phase set,
the same expression yields a fixed polyphase bank.

\subsection{Equivalence audit}

The binary64 audit uses 72 deterministic Gaussian three-channel lines over
lengths \(5,9,17,33,65,129\). The maximum difference from the formal current
is \(9.44\times10^{-16}\); the maximum synthesized-value difference is
\(8.88\times10^{-16}\). Cardinal error is zero, current-mass residual is at
most \(1.78\times10^{-15}\), and no sign-variation surplus is observed.

A second audit in Section~XI compares the local action compilation with the
formal all-source Eikonal operator. It reports the defect in
\eqref{eq:stardefect}, rather than conflating that approximation with the
binary64-equivalent FIR factor compilation.

\subsection{FIR refinement}

Collecting the preceding reductions, CONV$^{*}$ comprises compact FIR
analysis, exact segmented lineage, vectorized scalar-threshold admission,
polyphase synthesis, and the cardinal \(Q_1\) action coordinate
\eqref{eq:starcardinalcompilation}. Its one-dimensional factor is
algebraically identical to CONV; in two dimensions, the only departure is the
explicitly bounded defect \eqref{eq:stardefect}.

\subsection{Can CONV\texorpdfstring{$^{*}$}{*} Compete with EASU?}

EASU produces each target pixel in a single two-dimensional pass over a fixed
neighborhood \cite{lottes2021}. CONV$^{*}$, however, retains a barrier between
its Cartesian factors, since the second factor analyzes the nonlinear output
of the first. No throughput ordering follows from the formulas alone;
nevertheless, every nonlinear cell operation has now been reduced to the
fixed scalar-threshold and scan primitives specified above.

\subsubsection{Parallel lineage}

Sign inheritance is obtained by seeded prefix propagation, and the boundary
objective likewise has a scan form. For one transition window
\([p,q]\), let \(s^-\) and \(s^+\) be its left and right coarse signs and
define
\begin{align}
 P^-(b)&=\sum_{k=p}^{b-1}[\min(0,s^-a_k)]^2,\nonumber\\
 S^+(b)&=\sum_{k=b}^{q}[\min(0,s^+a_k)]^2.
 \label{eq:starconflictscan}
\end{align}
Using these partial sums, every candidate cost in \eqref{eq:boundary} becomes
\begin{equation}
 C(b)=P^-(b)+S^+(b),
 \label{eq:starboundarycost}
\end{equation}
so that one prefix sum and one suffix sum replace repeated evaluation of the
overlapping windows.

The dependence between successive transition boundaries is finite as well.
Indeed, if \(\kappa_\ell-\kappa_{\ell-1}\geq2\), then
\begin{equation}
 b_{\ell-1}+1\leq5\kappa_\ell-5<5\kappa_\ell-4,
 \label{eq:starwindowseparation}
\end{equation}
and the preceding boundary cannot restrict the candidate set
\eqref{eq:candidates}; dependence remains only inside a run of adjacent
coarse transitions. Writing
\(r_\ell=b_\ell-5\kappa_\ell\in\{-4,\ldots,4\}\), each transition in such a
run defines a deterministic map \(G_\ell\) from the preceding offset to the
next, including the stated first-minimum tie rule. Since function composition
is associative, a parallel prefix composition of these nine-state maps yields
the exact greedy result.

\subsubsection{Factor schedule}

The remaining work of one factor is fixed local analysis, the lineage scans,
the scalar solve \eqref{eq:thresholdmass}, and the polyphase dot product
\eqref{eq:starpolyphase}. Node jets may be formed once and shared by their two
incident cells. Fixed-scale phases are constants, and components in the
declared channel basis may be evaluated together. A tiled implementation may
write the synthesized factor through an on-chip transpose so that the next
factor again reads contiguous lines; it may not remove the intervening factor
barrier without changing the operator.

These reductions preserve the boundary closures, ordered tie rules, channel
basis, and Cartesian factor order of CONV. They reduce the remaining
nonlinear state to scans and a five-value piecewise-affine root. Whether that
arithmetic advantage offsets the intermediate factor traffic is a property
of a concrete implementation and target architecture. No performance claim
is made here; equations \eqref{eq:thresholdcurrent}--\eqref{eq:starwindowseparation}
specify the exact reduced form for such an implementation.
% END COMPOSED CONVSTAR

\section{Inversion in Conventional Resampling}
\enlargethispage{2\baselineskip}

Windowed-sinc, cubic, and edge-directed raster interfaces are ordinarily
specified as synthesis maps from one grid to another.\footnote{Glory first
be to The Holy One, Blessed be He
({\hebrewfont הַקָּדוֹשׁ בָּרוּךְ הוּא}), from whom all
good things come. To Him we attribute the ability that has provided for all
good to come and be demonstrated anywhere herein, and declare ourselves
unworthy servants. If in any way we have erred, or our work is insufficient,
we ask mercy and kindness, for our effort is good and earnest.} Their multiresolution
analysis and synthesis filters are separately constructed, and the identity
\(D_sU_s=I\) on coarse inputs does not make the discarded fine array
recoverable. CONV adds a native analysis state: conserved cell averages,
admitted centered moments, and residual moment fields have explicit analysis
and synthesis maps. The comparison is representational rather than a
claim that invertible multiscale methods are new; wavelet lifting already
establishes the general analysis--prediction--detail principle
\cite{sweldens1998}, and Harten's construction already represents nested
cell averages by coarse averages and fine-grid prediction errors
\cite{harten1995}. Here those established coordinates are joined to compact quintic
transport and a factorwise ringing-incapability theorem.

\section{Conclusion}

Compact Ordered Nodal Variation formulates interpolation and matched
multiresolution transfer through a conservative signed-current
representation. A high-order two-jet supplies the current proposal, ordered
admission restricts its derivative-sign lineage before Bernstein synthesis,
and a source-derived Riemannian coordinate couples the two Cartesian factor
orders. The same representation yields arbitrary-scale basin reduction,
exact residual-moment lifting, and the FIR/polyphase realization
CONV$^{*}$.

The analysis establishes cardinality, affine reproduction, cell-current
conservation, factorwise variation diminution, fixed-ledger nonexpansiveness,
sample-count-independent error bounds, well-posed Riemannian action, exact
retained-detail inversion, and Nyquist annihilation by conservative
restriction. The deterministic synthetic results confirm the predicted sign
topology and rapid resolved-field convergence; matched-cycle comparisons show
competitive fidelity with substantially smaller range excursion than
Lanczos on the selected diagonal interface, while native timings identify
the remaining cost of nonlinear admission and the Cartesian factor barrier.

Appendix~\ref{app:directwarp} develops the joint two-dimensional current
admission required for direct affine and projective warps, extending the
factorwise construction without intermediate rasterization.

% BEGIN COMPOSED DIRECT WARP APPENDIX
% Direct-warp appendix shared by the modular and composed paper builds.
\providecommand{\WarpFigureRoot}{conv_paper_assets}
\appendices
\section{Direct Joint CONV Warps}
\label{app:directwarp}

The continuous CONV representation admits a direct extension from lattice
resize to rotation, skew, anisotropic affine transfer, and perspective warp.
The necessary change is not a decomposition of the coordinate map into
one-dimensional resampling passes.  It is a joint admission of the
two-dimensional current before that map is evaluated.  The construction
below consequently retains one inverse map, one continuous source atlas, and
one positive raster functional.  No intermediate raster occurs.

\subsection{One inverse map and its induced geometry}

Let \(\Omega\subset\R^2\) denote the normalized source rectangle and let
\(\widehat\Omega\subset\R^2\) denote the target rectangle.  A nonsingular
inverse warp
\begin{equation}
 F:\widehat\Omega\longrightarrow\Omega,\qquad q\longmapsto p=F(q),
 \label{eq:warpmap}
\end{equation}
is declared once.  Given the admitted source atlas \(U_Y\), point synthesis is
\begin{equation}
 (\mathcal W_FY)(q)=U_Y(F(q)).
 \label{eq:warppoint}
\end{equation}
Thus a shear factorization of an affine map is neither required nor used.
Indeed, repeated admitted interpolation would generally define a different
nonlinear operator, because every intermediate raster changes the current
state from which the next admission is formed.

When the source current geometry is expressed by the SPD field \(M(p)\), the
same inverse map induces the target quadratic form
\begin{equation}
 G_F(q)=J_F(q)^TM(F(q))J_F(q),
 \qquad J_F=DF,
 \label{eq:warppullback}
\end{equation}
and its determinant-one shape is
\begin{equation}
 \overline M_F(q)=\frac{G_F(q)}{\det(G_F(q))^{1/2}}.
 \label{eq:warpshape}
\end{equation}
Equation~\eqref{eq:warppullback} combines sampling scale, anisotropy, and the
warp before a square root or directional support is formed.  Consequently it
does not assume that Cartesian spacing commutes with the eigenframe of
\(M\).  The scalar \(\det(G_F)^{1/4}\) carries linear scale, while
\eqref{eq:warpshape} carries directional shape.

\subsection{Compiled bivariate proposal}

Index the source cells by \(C\), and write \((u,v)\in[0,1]^2\) for local
cell coordinates.  The joint source potential on cell \(C\) is represented
in the tensor Bernstein basis as
\begin{equation}
 U_C(u,v)=\sum_{i=0}^{5}\sum_{j=0}^{5}
 P_{C,ji}B_i^5(u)B_j^5(v).
 \label{eq:warppatch}
\end{equation}
The proposal controls are not fitted by an optimizer.  Let
\begin{equation}
 A_{ri}=B_i^5(r/5),\qquad 0\leq r,i\leq5,
 \label{eq:warpcollocation}
\end{equation}
and let \(Z_C\) contain the values of the complete two-order CONV proposal at
the \(6\times6\) fifth-step tensor lattice of \(C\).  Since the Bernstein
collocation matrix \(A\) is nonsingular, the unique bidegree-\((5,5)\)
proposal is
\begin{equation}
 \widetilde P_C=A^{-1}Z_CA^{-T}.
 \label{eq:warpproposal}
\end{equation}
All entries of every \(Z_C\) are obtained from one global fifth-step
evaluation.  Adjacent cells therefore have the same six samples on their
common edge; uniqueness in \eqref{eq:warpcollocation} then gives the same
complete Bernstein edge.  Equivalently, the patches are stored in one shared
global control lattice.  The resulting proposal is \(C^0\), and its current
is the gradient of one continuous piecewise-polynomial potential.

Let \(B\) denote the degree-elevated \(Q_1\) interpolant of the source
samples on that shared lattice.  It is cardinal, continuous, and affine
exact.  Admission will move \(B\) toward \(\widetilde P\), while the source
vertices remain fixed.

\subsection{Support-range admission}

For each source cell \(C\), let \(\mathcal S_C\) be the complete Cartesian
two-jet sample support, clipped at the outer boundary.  If shared control
entity \(e\) is incident to the cell set \(\mathcal I(e)\), define, for
channel \(c\),
\begin{equation}
 I_e^{(c)}=
 \bigcap_{C\in\mathcal I(e)}
 \left[
  \min_{p\in\mathcal S_C}Y_p^{(c)},
  \max_{p\in\mathcal S_C}Y_p^{(c)}
 \right].
 \label{eq:warprangeintersection}
\end{equation}
This intersection is nonempty: the corresponding \(Q_1\) control is a
convex combination of common incident vertices and belongs to every interval
in \eqref{eq:warprangeintersection}.  The proposal is first replaced by
\begin{equation}
 \widehat P_e^{(c)}=\Pi_{I_e^{(c)}}\widetilde P_e^{(c)}.
 \label{eq:warprangeclip}
\end{equation}
Unlike a post-reconstruction clamp, \eqref{eq:warprangeclip} constrains the
continuous potential.  Subsequent current admission remains on the segment
from \(B_e^{(c)}\) to \(\widehat P_e^{(c)}\), so every final control remains
in \(I_e^{(c)}\).

\begin{theorem}[Continuous support-range admission]
For every channel and source cell,
\begin{equation}
 \min_{p\in\mathcal S_C}Y_p^{(c)}
 \leq U_C^{(c)}(u,v)\leq
 \max_{p\in\mathcal S_C}Y_p^{(c)},
 \qquad (u,v)\in[0,1]^2.
 \label{eq:warprangetheorem}
\end{equation}
Consequently, point synthesis and every positive warped-pixel average have no
range excursion beyond the union of the source supports they traverse.
\end{theorem}

\begin{proof}
For cell \(C\), every one of its 36 shared controls belongs to the interval
on the right of \eqref{eq:warprangetheorem}.  The tensor Bernstein basis is
nonnegative and forms a partition of unity, hence \(U_C\) is a convex
combination of those controls.  Point evaluation gives the first result, and
integration against a nonnegative normalized measure gives the second.
\end{proof}

\subsection{Support-current cone}

The one-dimensional theorem controls current along a declared factor.  A
rotated or projective query instead requires a joint source current law.  For
cell \(C\), collect the four corner gradients of every \(Q_1\) cell contained
in \(\mathcal S_C\):
\begin{equation}
 \mathcal G_C=
 \left\{\nabla Q_D(z):D\subset\mathcal S_C,
              \ z\in\operatorname{vert}(D)\right\}.
 \label{eq:warpgenerators}
\end{equation}
Zero vectors are removed.  If the remaining directions generate a proper
closed planar cone, let
\begin{equation}
 K_C=\operatorname{cone}\mathcal G_C;
 \label{eq:warpcone}
\end{equation}
if no nonzero current is observed, or if the observed currents positively
span the plane, set \(K_C=\R^2\).  Thus a clean interface gives a narrow
cone, curvature gives a wedge, and crossing currents remove a directional
restriction rather than forcing a false owner or normal.

Every proper planar cone occurring here has a finite half-space description
\begin{equation}
 K_C=\left\{z\in\R^2:\ell_{Cr}^Tz\geq0,
                       \ 1\leq r\leq R_C\right\},
 \qquad R_C\leq3.
 \label{eq:warpconedual}
\end{equation}
A wedge uses its two inward boundary normals.  A ray additionally requires a
forward inequality, while a line uses two opposing transverse inequalities.
The whole-plane case has \(R_C=0\).

For each \(r\), the directional derivative
\begin{equation}
 L_{Cr}(u,v)=\ell_{Cr}^T\nabla U_C(u,v)
 \label{eq:warpdirectional}
\end{equation}
is degree elevated to the bidegree-\((5,5)\) Bernstein basis.  Denote its 36
linear control functionals by \(a_{Crk}^TP\).  The finite admission
conditions are
\begin{equation}
 a_{Crk}^TP\geq0,\qquad
 1\leq r\leq R_C,\quad 1\leq k\leq36.
 \label{eq:warpfiniteconstraints}
\end{equation}
Nonnegativity of these Bernstein coefficients is a sufficient certificate
for \(L_{Cr}\geq0\) on the complete cell.

The degree-elevated \(Q_1\) baseline is a constructive feasible point.  Its
gradient is a convex combination of the four corner gradients of its source
cell; those gradients occur in \(\mathcal G_C\).  Hence
\begin{equation}
 \nabla B_C(u,v)\in K_C
 \label{eq:warpbaselinefeasible}
\end{equation}
everywhere, and every margin in \eqref{eq:warpfiniteconstraints} is
nonnegative.

\subsection{Finite topological-entity admission}

Assign each nonvertex control to exactly one shared-edge or cell-interior
entity \(g\).  With the range-admitted proposal from
\eqref{eq:warprangeclip}, write
\begin{equation}
 \widehat P=B+\sum_gd_g,
 \label{eq:warpgroups}
\end{equation}
where the entity corrections have disjoint control ownership.  For a global
constraint \(a_k^TP\geq0\), put
\begin{equation}
 m_k=a_k^TB,\qquad
 H_k=\sum_g\max(0,-a_k^Td_g).
 \label{eq:warpmargin}
\end{equation}
If \(\widehat P\) already satisfies all constraints, it is accepted
unchanged.  Otherwise each harmful entity is assigned the simultaneous
capacity
\begin{equation}
 \alpha_g=min\!\left\{1,
  \min_{k:a_k^Td_g<0,\ H_k>0}\frac{m_k}{H_k}\right\},
 \label{eq:warpcapacity}
\end{equation}
with an empty inner minimum interpreted as one.  The first admitted state is
\begin{equation}
 P^{(0)}=B+\sum_g\alpha_gd_g.
 \label{eq:warpfirststate}
\end{equation}
Finally, let \(R=\widehat P-P^{(0)}\) and take the maximal common ray
\begin{equation}
 \tau=\min\!\left\{1,
   \min_{k:a_k^TR<0}
   \frac{a_k^TP^{(0)}}{-a_k^TR}\right\},
 \qquad P=P^{(0)}+\tau R.
 \label{eq:warpray}
\end{equation}
Again, an empty inner minimum is one.  Equations
\eqref{eq:warpcapacity}--\eqref{eq:warpray} contain neither an iteration nor
an active-set search.

\begin{theorem}[Finite joint current admission]
The control lattice \(P\) obtained from
\eqref{eq:warpcapacity}--\eqref{eq:warpray} satisfies every constraint in
\eqref{eq:warpfiniteconstraints}.  Consequently,
\begin{equation}
 \nabla U_C(u,v)\in K_C,
 \qquad (u,v)\in[0,1]^2,
 \label{eq:warpgradientcone}
\end{equation}
for every source cell.  The construction is finite, order independent,
cardinal, \(C^0\), and affine exact.
\end{theorem}

\begin{proof}
Fix constraint \(k\).  Contributions with \(a_k^Td_g\geq0\) can be
discarded in a lower bound.  For every harmful entity,
\(\alpha_g\leq m_k/H_k\), and therefore
\begin{equation}
 a_k^TP^{(0)}
 \geq m_k-\frac{m_k}{H_k}
       \sum_g\max(0,-a_k^Td_g)
 =0.
 \label{eq:warpcapacityproof}
\end{equation}
Equation~\eqref{eq:warpray} is the exact first intersection of the remaining
affine ray with the finite half-space boundary, hence \(a_k^TP\geq0\) as
well.  Applied to all derivative coefficients, Bernstein positivity gives
\(\ell_{Cr}^T\nabla U_C\geq0\), which is equivalent to
\eqref{eq:warpgradientcone} through \eqref{eq:warpconedual}.

Source vertices are not assigned to a movable entity.  Shared-edge controls
are single variables, so cardinality and \(C^0\) continuity are retained.
For affine data, the proposal equals the degree-elevated affine baseline and
already satisfies the cone and range constraints; it is consequently passed
unchanged.
\end{proof}

The theorem also explains why the construction is not an edge detector.  It
does not select one current from \(\mathcal G_C\).  Instead it admits the
complete high-order proposal whenever all of its derivative coefficients lie
in the cone generated by the currents actually present in the declared
support.

\subsection{Transport through affine and projective maps}

Let \(\gamma:[a,b]\to\widehat\Omega\) be a \(C^1\) target curve.  At every
regular point contained in one source cell, the chain rule gives
\begin{equation}
 \frac{d}{ds}U_Y(F(\gamma(s)))=
 \bigl[J_F(\gamma(s))\gamma'(s)\bigr]^T
 \nabla U_Y(F(\gamma(s))).
 \label{eq:warpchain}
\end{equation}
Let
\begin{equation}
 K_C^*=\{t\in\R^2:t^Tz\geq0\ \text{for all }z\in K_C\}
 \label{eq:warpdual}
\end{equation}
be the dual current direction cone.

\begin{corollary}[Warped-current transport]
Suppose \(F(\gamma(s))\) crosses only finitely many source-cell boundaries
and
\begin{equation}
 J_F(\gamma(s))\gamma'(s)\in K_{C(s)}^*
 \label{eq:warptangentcondition}
\end{equation}
for almost every \(s\), where \(C(s)\) contains \(F(\gamma(s))\).  Then
\(U_Y\circ F\circ\gamma\) is nondecreasing on \([a,b]\).
\end{corollary}

\begin{proof}
Equations \eqref{eq:warpgradientcone}, \eqref{eq:warpdual}, and
\eqref{eq:warpchain} make the derivative nonnegative in every open patch
segment.  The atlas is continuous and piecewise \(C^1\); therefore its
composition with the curve is absolutely continuous after splitting at the
finite knot set.  The fundamental theorem gives
\begin{equation}
 U_Y(F(\gamma(b)))-U_Y(F(\gamma(a)))
 =\int_a^b\frac{d}{ds}U_Y(F(\gamma(s)))\,ds\geq0.
\end{equation}
The knot points have measure zero and require no derivative convention.
\end{proof}

Affine maps transport the admitted directions by one constant Jacobian.
Homographies transport them by the spatially varying Jacobian of the same
declared map.  In either case, no target-space owner, orientation, or second
interpolation operation is introduced.

\subsection{Conservative warped-pixel analysis}

Point synthesis is appropriate for enlargement.  For reduction, let
\(\{P_j\}\) be the endpoint-aligned target Voronoi basins.  The matched
raster functional is
\begin{equation}
 (\mathcal B_FY)_j=
 \frac{1}{|P_j|}\int_{P_j}U_Y(F(q))\,dq.
 \label{eq:warppixel}
\end{equation}
Since the basins partition the target domain,
\begin{equation}
 \sum_j|P_j|(\mathcal B_FY)_j
 =\int_{\cup_jP_j}U_Y(F(q))\,dq.
 \label{eq:warpconservation}
\end{equation}
Thus \eqref{eq:warppixel} is conservative relative to the admitted continuous
profile.  It is also a positive average and inherits
\eqref{eq:warprangetheorem}.

For affine \(F\), a target basin maps to a source parallelogram.  Splitting
at the integer source-knot lines produces convex polygons, each contained in
one patch.  Positive fan triangulation introduces no overlap.  On every
triangle, \(U_C\) has total degree at most ten.  Under a Duffy map the area
Jacobian raises one separate degree to at most eleven; consequently, the
tensor three-point rule is insufficient, whereas tensor Gauss--Legendre
order six is exact.  Division by the mapped parallelogram area then gives
\eqref{eq:warppixel} exactly.

For a homography, source-coordinate integration yields
\begin{equation}
 (\mathcal B_FY)_j=\frac1{|P_j|}
 \int_{F(P_j)}U_Y(p)\left|\det DF^{-1}(p)\right|\,dp.
 \label{eq:warpprojectiveintegral}
\end{equation}
With homogeneous coordinate normalizations absorbed into the nonsingular
matrix \(H\), the inverse-map density is
\begin{equation}
 \left|\det DF^{-1}(p)\right|
 =\frac{|\det H^{-1}|}{|\ell(p)|^3},
 \label{eq:warpdensity}
\end{equation}
where \(\ell\) is the affine denominator of \(H^{-1}\).  On a clipped
triangle choose a nonzero center \(c\) of the denominator interval and write
\(\ell=c(1+z)\), with \(|z|\leq r<1\).  Then
\begin{equation}
 (1+z)^{-3}=\sum_{k=0}^{K}(-1)^k\binom{k+2}{2}z^k+R_K(z).
 \label{eq:warpseries}
\end{equation}
The tail admits the closed bound
\begin{align}
 |R_K(z)|&\leq T_{K+1}(r),\label{eq:warptail}\\
 T_n(r)&=\frac{r^n}{2}\left[
 \frac{r(1+r)}{(1-r)^3}
 +\frac{(2n+3)r}{(1-r)^2}\right.\nonumber\\
 &\hspace{25mm}\left.
 +\frac{(n+1)(n+2)}{1-r}\right].\nonumber
\end{align}
The finite series multiplied by \(U_C\) is polynomial and is integrated
exactly by tensor Gauss order
\begin{equation}
 q_K=\left\lceil\frac{K+12}{2}\right\rceil.
 \label{eq:warpseriesorder}
\end{equation}
If \(U_C=P_0+(U_C-P_0)\), the constant \(P_0\) is integrated exactly before
\eqref{eq:warpseries} is applied.  Hence constants have zero certified
residual.  For triangle \(T\), the remaining error is bounded by
\begin{equation}
 E_T\leq |T|\frac{|\det H^{-1}|}{|c|^3}
 T_{K+1}(r)\,
 \norm{U_C-P_0}_{L^\infty(T)}.
 \label{eq:warpprojectivecertificate}
\end{equation}
Bernstein controls and first differences provide the required value bound.
Four-way triangle subdivision is invoked only when the sum of
\eqref{eq:warpprojectivecertificate} exceeds the declared pixel budget.  The
exact integral \eqref{eq:warpprojectiveintegral} remains the operator; the
finite implementation returns its value together with a rigorous enclosure.

\subsection{Multichannel form}

For \(Y:\Omega\cap\mathbb Z^2\to\R^C\), the control topology, support sets,
and inverse map are shared.  Range and current admission are applied to each
component in the declared channel basis, yielding the product set
\begin{equation}
 \prod_{c=1}^C
 \left\{U^{(c)}:\nabla U_C^{(c)}\in K_C^{(c)},\ 
 U_C^{(c)}\in I_C^{(c)}\right\}.
 \label{eq:warpmultichannel}
\end{equation}
The preceding proofs therefore hold componentwise.  As with the main CONV
operator, this statement is invariant under independent affine changes of
channel units; it does not assert invariance under an arbitrary basis change
that mixes channels.

\subsection{Exact-truth evaluation}

The evidence uses only procedural sources whose point values and warped-pixel
means are independent of every tested interpolator.  The sharp corpus
contains a half-plane with normal angle \(37^\circ\), a width-0.065 bar at
\(27^\circ\), two width-0.055 crossing bars at \(31^\circ\) and
\(121^\circ\), and a \(0.62\times0.34\) box at \(-19^\circ\), all centered
at \((1/2,1/2)\).  The four target-to-source maps are
\begin{align}
 F_{\rm rot}&: (\theta,s_x,s_y)=(21^\circ,0.72,0.72),\nonumber\\
 F_{\rm skew}&: (s_x,s_y,h_x)=(0.70,0.75,0.42),\nonumber\\
 F_{\rm aniso}&: (\theta,s_x,s_y)=(-14^\circ,0.50,0.84),\nonumber\\
 F_{\rm proj}&: (g_x,g_y,\theta,s_x,s_y)
 =(0.32,-0.18,9^\circ,0.67,0.72),
 \label{eq:warpcorpusmaps}
\end{align}
where affine maps use \(F(q)=\tfrac12+A(q-\tfrac12)\), with
\(A=R_\theta\left(\begin{smallmatrix}1&h_x\\h_y&1\end{smallmatrix}\right)
\operatorname{diag}(s_x,s_y)\), and the projective map uses
\begin{equation}
 F(q)=\frac12+
 \frac{A(q-\frac12)}{1+(g_x,g_y)^T(q-\frac12)}.
 \label{eq:warpprojectivedeclared}
\end{equation}
Exact polygon membership supplies point truth.  Since a source half-plane
remains a target half-plane under a pole-free homography, polygon clipping in
target coordinates supplies exact basin coverage.

The smooth census uses the real Fourier fields
\begin{equation}
 f(x)=b+\sum_{k=1}^Ka_k
 \cos\!\left(2\pi\xi_k^Tx+\varphi_k\right).
 \label{eq:warpfourierfield}
\end{equation}
For an affine map, the exact mean over a target rectangle with center \(m\)
and side lengths \(h_x,h_y\) is obtained by multiplying each transformed
mode by
\begin{equation}
 \operatorname{sinc}((A^T\xi_k)_xh_x)
 \operatorname{sinc}((A^T\xi_k)_yh_y).
 \label{eq:warpfouriermean}
\end{equation}
The scalar probes use the two mode sets
\begin{equation}
\begin{aligned}
 \mathcal F_1=\{&((5,2),.22,.13),\ ((9,3),.09,1.17),\\
                &((2,-4),.07,-.41)\},\\
 \mathcal F_2=\{&((6,1),.17,.2),\ ((-1,6),.17,-.6),\\
                &((4,4),.09,1.1)\},
\end{aligned}
\label{eq:warpfourierprobes}
\end{equation}
with \(b=0.5\).  A three-channel probe uses wavevectors
\((4,1),(1,5),(5,-3)\), phases \(0,0.7,-0.4\), and amplitude rows
\begin{equation}
 \begin{bmatrix}.22&.05&.08\\.04&.20&.07\\.06&.04&.18\end{bmatrix},
 \qquad b=(.5,.5,.5).
 \label{eq:warpmultichannelprobe}
\end{equation}

Table~\ref{tab:warpsharp} reports the 16 sharp cases on a
\(17\times17\) source and \(13\times13\) target.  CONV has the lowest MSE in
seven point cases and seven warped-pixel cases.  More importantly for the
proved admission, its maximum range excursion remains at floating roundoff,
whereas the factor-only ablation and Lanczos admit substantial excursion.

\begin{table}[H]
\caption{Sharp affine/projective exact-truth census}
\label{tab:warpsharp}
\centering
\scriptsize
\setlength{\tabcolsep}{2.0pt}
\begin{tabular}{lrrrr}
\hline
Method & Point MSE & Pixel MSE & Point exc. & Pixel exc.\\
\hline
CONV & .021002 & .007377 & \(2.22\!\times\!10^{-16}\) & \(3.18\!\times\!10^{-14}\)\\
Factor ablation & .021137 & .007705 & .2287 & .1433\\
Lanczos-3 & .022928 & .008741 & .2817 & .1989\\
Bilinear & .021513 & .007576 & 0 & 0\\
\hline
\end{tabular}
\end{table}

\begin{figure*}[!t]
 \centering
 \includegraphics[width=0.98\textwidth]{conv_paper_assets/conv_warp_projective_plate_canonical.png}
 \vspace{3pt}
 \includegraphics[width=0.74\textwidth]{conv_paper_assets/conv_warp_canonical_metrics.png}
 \caption{Direct joint CONV warp evidence.  Top: one homography applied to
 sources with exact projective basin truth; each displayed certificate is the
 computed enclosure of the projective polynomial integral.  Bottom:
 topology and exact-truth fidelity over the declared corpus.  Values plotted
 at the numerical floor are below \(2\times10^{-16}\).  The factor-only
 result is an ablation of joint current admission, not a second CONV
 definition.}
 \label{fig:warpevidence}
\end{figure*}

The intrinsic topology census evaluates normal profiles rather than using MSE
as a surrogate for ringing.  For a profile \(z_0,\ldots,z_n\) whose exact
orientation is increasing, the reported wrong-way current is
\begin{equation}
 V_-(z)=\sum_{i=0}^{n-1}\max(0,z_i-z_{i+1}).
 \label{eq:warpwrongway}
\end{equation}
Straight interfaces use six source angles, three maps, 513 samples per
profile, and source sides 17 and 33.  CONV has no normal sign changes,
\(V_-\leq1.34\times10^{-14}\), and excursion below
\(4.45\times10^{-16}\).  Lanczos-3 produces four to nine sign changes,
mean \(V_-=0.276\)--\(0.277\), and mean excursion \(0.166\).

For curved interfaces, circles and ellipses are tested at 12 phases under
three maps, giving 72 profiles per source size.  Table~\ref{tab:warpcurved}
shows that the admitted cone eliminates wrong-way normal current at sides 13,
17, and 25.  At side 9, the range theorem remains exact, while the six-node
support spans competing conic normals and the cone correctly widens; the
remaining wrong-way current therefore measures support ambiguity rather than
failure of the finite inequalities.

\begin{table}[H]
\caption{Curved-interface intrinsic normal census for CONV}
\label{tab:warpcurved}
\centering
\footnotesize
\setlength{\tabcolsep}{4pt}
\begin{tabular}{rrrr}
\hline
Source side & Mean \(V_-\) & Max. \(V_-\) & Max. excursion\\
\hline
9  & \(1.94\times10^{-3}\) & \(2.80\times10^{-2}\) & \(2.22\times10^{-16}\)\\
13 & \(1.67\times10^{-15}\) & \(4.77\times10^{-15}\) & \(4.44\times10^{-16}\)\\
17 & \(2.23\times10^{-15}\) & \(6.44\times10^{-15}\) & \(4.44\times10^{-16}\)\\
25 & \(3.31\times10^{-15}\) & \(5.44\times10^{-15}\) & \(4.44\times10^{-16}\)\\
\hline
\end{tabular}
\end{table}

On the nine band-admitted Fourier cases, mean point MSE is
\(1.48764\times10^{-4}\) for CONV, \(1.06820\times10^{-4}\) for the
factor-only ablation, \(4.57424\times10^{-5}\) for Lanczos-3, and
\(1.22893\times10^{-3}\) for bilinear.  The corresponding warped-pixel MSEs
are \(6.27668\times10^{-5}\), \(4.09234\times10^{-5}\),
\(1.88691\times10^{-5}\), and \(7.63065\times10^{-4}\).  Thus the complete
support maximum principle has a measurable smooth-field fidelity cost; the
experiment does not support universal MSE dominance.

Finally, affine conservation equates the weighted target sum and an
independent exact integral of the admitted atlas to \(3\times10^{-12}\).  On
a nonconstant projective plane, independent 40-point tensor integration
differs from the certified result by \(1.18\times10^{-11}\), within its
\(5.84\times10^{-10}\) enclosure.  The largest certificate returned by the
sharp projective corpus is \(1.52\times10^{-11}\).  These computations verify
the exact polynomial and error-enclosure implementations; the cardinality,
integrability, range, current-cone, and transport statements follow from the
preceding formulas rather than from the measurements.

The direct warp therefore extends CONV without changing its governing idea:
the high-order proposal is admitted in the representation that carries the
current, after which point evaluation and positive basin analysis are two
functionals of the same continuous object.  The principal remaining boundary
is source resolution itself.  When one compact two-jet support contains
genuinely competing interface normals, the observed cone contains both and
cannot certify a narrower direction without adding information absent from
the samples.
% END COMPOSED DIRECT WARP APPENDIX

\section*{AI GENERATED WORK DISCLAIMER}

Substantial portions of this work were generated with the assistance of
OpenAI Codex general reasoning models.

\begin{thebibliography}{99}
\scriptsize
\bibitem{shannon1949}
C. E. Shannon, "Communication in the presence of noise," \emph{Proc. IRE},
vol. 37, no. 1, pp. 10--21, Jan. 1949.
\bibitem{lanczos1956}
C. Lanczos, \emph{Applied Analysis}. Englewood Cliffs, NJ, USA:
Prentice-Hall, 1956.
\bibitem{keys1981}
R. G. Keys, "Cubic convolution interpolation for digital image
processing," \emph{IEEE Trans. Acoust., Speech, Signal Process.}, vol. 29,
no. 6, pp. 1153--1160, Dec. 1981.
\bibitem{fritsch1980}
F. N. Fritsch and R. E. Carlson, "Monotone piecewise cubic interpolation,"
\emph{SIAM J. Numer. Anal.}, vol. 17, no. 2, pp. 238--246, Apr. 1980.
\bibitem{hyman1983}
J. M. Hyman, "Accurate monotonicity preserving cubic interpolation,"
\emph{SIAM J. Sci. Stat. Comput.}, vol. 4, no. 4, pp. 645--654, Dec. 1983.
\bibitem{costantini1987}
P. Costantini, "Co-monotone interpolating splines of arbitrary degree---a
local approach," \emph{SIAM J. Sci. Stat. Comput.}, vol. 8, no. 6,
pp. 1026--1034, Nov. 1987, doi: 10.1137/0908083.
\bibitem{dougherty1989}
R. Dougherty, A. Edelman, and J. M. Hyman, "Nonnegativity-,
monotonicity-, or convexity-preserving cubic and quintic Hermite
interpolation," \emph{Math. Comp.}, vol. 52, no. 186, pp. 471--494,
Apr. 1989, doi: 10.1090/S0025-5718-1989-0962209-1.
\bibitem{ulrich1994}
G. Ulrich and L. T. Watson, "Positivity conditions for quartic
polynomials," \emph{SIAM J. Sci. Comput.}, vol. 15, no. 3, pp. 528--544,
May 1994, doi: 10.1137/0915035.
\bibitem{lux2023}
T. Lux, L. T. Watson, T. Chang, and W. Thacker, "Algorithm 1031:
MQSI---monotone quintic spline interpolation," \emph{ACM Trans. Math.
Softw.}, vol. 49, no. 1, Art. no. 6, pp. 1--17, Mar. 2023,
doi: 10.1145/3570157.
\bibitem{carlsonfritsch1989}
R. E. Carlson and F. N. Fritsch, "An algorithm for monotone piecewise
bicubic interpolation," \emph{SIAM J. Numer. Anal.}, vol. 26, no. 1,
pp. 230--238, Feb. 1989, doi: 10.1137/0726013.
\bibitem{costantinimanni1991}
P. Costantini and C. Manni, "A local scheme for bivariate co-monotone
interpolation," \emph{Comput. Aided Geom. Des.}, vol. 8, no. 5,
pp. 371--391, Nov. 1991, doi: 10.1016/0167-8396(91)90011-Y.
\bibitem{li2001}
X. Li and M. T. Orchard, "New edge-directed interpolation," \emph{IEEE
Trans. Image Process.}, vol. 10, no. 10, pp. 1521--1527, Oct. 2001.
\bibitem{lottes2021}
T. Lottes, "AMD FidelityFX Super Resolution 1.0," Advanced Micro Devices,
Santa Clara, CA, USA, 2021.
\bibitem{ding2007}
W. Ding, F. Wu, X. Wu, S. Li, and H. Li, "Adaptive directional
lifting-based wavelet transform for image coding," \emph{IEEE Trans. Image
Process.}, vol. 16, no. 2, pp. 416--427, Feb. 2007,
doi: 10.1109/TIP.2006.888341.
\bibitem{bigun1987}
J. Bigun and G. H. Granlund, ``Optimal orientation detection of linear
symmetry,'' in \emph{Proc. 1st IEEE Int. Conf. Comput. Vis.}, London, U.K.,
1987, pp. 433--438.
\bibitem{mirebeau2014}
J.-M. Mirebeau, ``Anisotropic fast-marching on Cartesian grids using lattice
basis reduction,'' \emph{SIAM J. Numer. Anal.}, vol. 52, no. 4,
pp. 1573--1599, 2014, doi: 10.1137/120861667.
\bibitem{shepard1968}
D. Shepard, ``A two-dimensional interpolation function for irregularly-spaced
data,'' in \emph{Proc. 23rd ACM Natl. Conf.}, New York, NY, USA, 1968,
pp. 517--524, doi: 10.1145/800186.810616.
\bibitem{karlin1968}
S. Karlin, \emph{Total Positivity}, vol. 1. Stanford, CA, USA: Stanford
Univ. Press, 1968.
\bibitem{pinkus2010}
A. Pinkus, \emph{Totally Positive Matrices}. Cambridge, U.K.: Cambridge
Univ. Press, 2010.
\bibitem{cavaretta1987}
A. S. Cavaretta, Jr. and A. Sharma, "Variation diminishing properties and
convexity for the tensor product Bernstein operator," in \emph{Functional
Analysis and Operator Theory}, Lecture Notes in Mathematics, vol. 1511.
Berlin, Germany: Springer, 1992, pp. 18--32,
doi: 10.1007/BFb0093794.
\bibitem{goodman1987}
T. N. T. Goodman, "Variation diminishing properties of Bernstein
polynomials on triangles," \emph{J. Approx. Theory}, vol. 50, no. 2,
pp. 111--126, 1987, doi: 10.1016/0021-9045(87)90002-5.
\bibitem{harten1995}
A. Harten, "Multiresolution algorithms for the numerical solution of
hyperbolic conservation laws," \emph{Commun. Pure Appl. Math.}, vol. 48,
no. 12, pp. 1305--1342, Dec. 1995,
doi: 10.1002/cpa.3160481201.
\bibitem{sweldens1998}
W. Sweldens, "The lifting scheme: A construction of second generation
wavelets," \emph{SIAM J. Math. Anal.}, vol. 29, no. 2, pp. 511--546, 1998,
doi: 10.1137/S0036141095289051.
\bibitem{gilbargtrudinger2001}
D. Gilbarg and N. S. Trudinger, \emph{Elliptic Partial Differential Equations
of Second Order}. Berlin, Germany: Springer, 2001.
\end{thebibliography}
\end{document}
